\documentclass[10pt, twoside, a4paper]{article}
\usepackage[italian]{babel}
\usepackage[utf8]{inputenc}
\usepackage{amsmath}
\usepackage{amsfonts}
\usepackage{fullpage}
\usepackage[pdftex]{graphicx}
\usepackage{booktabs}
\usepackage{wrapfig}
\usepackage{multirow}
\usepackage{sidecap}
\sidecaptionvpos{figure}{t}
\usepackage{subcaption}
\usepackage{siunitx}
\usepackage[font=small]{caption}
\usepackage[bookmarks, hidelinks]{hyperref}
\usepackage{float}
%nuovi pacchetti
\usepackage{fontenc}
\usepackage{fancyhdr}
\usepackage{amssymb}
\usepackage{enumitem}
\usepackage[a4paper, top=1.4cm, bottom=1.4cm, left=1.5cm, right=1.5cm] {geometry}
\usepackage{multicol}
%pacchetti aggiunti da poco
\usepackage{footnote}

% COMANDI PER PREMETTERE AL NUMERO DELLE FORMULE E DELLE FIGURE IL NUMERO DELLA SEZIONE
\numberwithin{equation}{section}
\numberwithin{figure}{section}
\numberwithin{table}{section}

% COMANDO PER SETTARE IN AUTOMATICO LE PARENTESI E IL \pm in \SI{num(inc)ESP}{\unit}
% esempio (10.00 \pm 0.05) m --> \SI{10.00(5)E0}{\metre}
\sisetup{separate-uncertainty,bracket-numbers = false}

\begin{document}

\title{Quaderno di laboratorio}
\date{A.A. 2014 - 2015}
\author{Davide Bazzanella}

\begin{titlepage}
\begin{center}

	\hrule \vspace{0.5cm}
     	\textsc{\LARGE QUADERNO DI LABORATORIO}
	\vspace{0.5cm} \hrule \vspace{2cm}

      	{\large Davide Bazzanella\\
		Gruppo A10}\\
	\vspace{0.5cm}
      	{\large A.A. 2014 - 2015}
	\vfill

	\includegraphics[width=4cm]{unitn_logo.png}\\
	\vspace{1cm}
        \textsc{\Large Università degli studi di Trento}
	\vfill

	%{\begin{abstract}
%Verifica della legge di stato dei gas ideali per trasformazioni isocore.

%Estrapolazione del valore dello zero assoluto con i dati ottenuti.
%	 \end{abstract}}
\end{center}
\end{titlepage}

\newpage
\input{tableofcontents.tex}
\newpage

\section{Amplificatori Operazionali Ideali - 16.09.2014}

In questa prima sessione di laboratorio abbiamo preparato la breadboard con il circuito di alimentazione per l'amplificatore operazionale (opamp) che useremo anche nelle successive esperienze.
Una volta preparata l'alimentazione, abbiamo montato due circuiti con amplificatori operazionali e valutato il loro funzionamento misurando la loro tensione di output.

\subsection{Strumenti e materiali}

\begin{itemize} [noitemsep]
\item Oscilloscopio Agilent DSO-X 2002A (bandwidth \SI{70}{\mega\hertz}, sample rate \num{2} GSa/s);
\item Generatore di tensione continua Agilent E3631A (max $\pm \, \SI{25}{\volt}$ o $\pm \, \SI{6}{\volt}$);
\item Generatore di forme d'onta Agilent 33120A con range di frequenza da \SI{100}{\micro\hertz} a \SI{15}{\mega\hertz};
\item Multimetro Agilent 34410A (utilizzato come amperometro e per verificare i valori delle resistenze);
\item Un amplificatore operazionale $\mu$A741;
\item Resistenze di vari valori;
\item Due capacità da \SI{100}{\nano\farad};
\item un trimmer (potenziometro);
\item Breadboard e cablaggi vari.
\end{itemize}

\subsection{Premessa sugli amplificatori operazionali ideali}

Nelle nostre esperienze tratteremo gli amplificatori operazionali considerandoli componenti ideali. Tale approssimazione non è limitante visti i valori di corrente e gli errori in gioco nel nostro caso.
Nel modello ideale il guadagno è infinito, così come anche l'impedenza di ingresso, mentre l'impedenza di uscita è nulla e pure il guadagno di modo comune.
Infine per un amplificatore operazionale ideale provvisto di circuito di retroazione negativo valgono le seguenti \textit{regole d'oro} (siano A e B rispettivamente gli ingressi invertente e non invertente):

\begin{equation}
\Delta V_{AB}=0
\label{eq:regola_V}
\end{equation}
\begin{equation}
I_{AB}=0
\label{eq:regola_I}
\end{equation}

Di queste due equazioni, la (\ref{eq:regola_V}) asserisce che la differenza di potenziale fra l'ingresso invertente e il non invertente è resa nulla dall'amplificatore tramite il circuito di retroazione negativo.
La seconda equazione, invece, afferma che la corrente assorbita dai terminali di ingresso dell'amplificatore è nulla.

Queste \textit{regole d'oro} sono indispensabili per l'analisi della risposta in uscita del circuito ad un segnale in entrata e verranno utilizzate ogniqualvolta verrà applicato il modello di opamp ideale.

\begin{SCfigure}[][ht]
 \centering
	\includegraphics[width=6.5cm]{../E01/latex/alimentazione.pdf}
	\caption{In questo grafico è mostrato lo schema di alimentazione dell'opamp. Le capacità hanno un valore nominale di $C_1=C_2= \SI{100}{\nano\farad}$. La tensione di alimentazione è impostata a $V^+ = \SI{15}{\volt}$ e a $V^- = \SI{-15}{\volt}$ ed è fornita dal generatore di tensione costante.\\
\newline\newline\newline}
 \label{gr:costante}
\end{SCfigure}

Per quanto riguarda l'alimentazione del circuito, abbiamo collegato i terminali $V^+$ e $V^-$ ad alimentazioni rispettivamente di \SI{+15}{\volt} e \SI{-15}{\volt}. Inoltre, in quanto immersi in un ambiente elettronicamente rumoroso, abbiamo inserito nel circuito due capacitori con lo scopo di evitare problemi di noise elettronico.
Per maggiore chiarezza di esposizione, questa configurazione sarà nascosta negli schemi successivi, ma comunque presente sulla breadboard.

%\newpage
\subsection{Generatore di corrente}

Come primo circuito abbiamo montato un generatore di corrente costante, ossia un dispositivo in grado di erogare una corrente costante indipendentemente dal carico a cui è collegato.
Poiché il nostro interesse era osservare il comportamento del circuito per diversi valori della resistenza di carico $R_2$, abbiamo utilizzato \textit{trimmer} come resistenza variabile. In Figura \ref{fig:gen_continua} è riportato lo schema circuitale.

\begin{wrapfigure}[18]{r}{0.52\textwidth}
  \begin{center}
    \includegraphics[width=0.42\textwidth]{../E01/latex/c1.pdf}
  \end{center}
  \caption{Schema del generatore di corrente costante. Come valori abbiamo utilizzato $R1=3.85 \pm 0.01 \,\si{\kilo\ohm}$ e $V_{gen}=3.85 V$, mentre $R_2$ è variabile. Come amperometro è stato utilizzato il multimetro, mentre sia per alimentare l'opamp sia come generatore di tensione costante, abbiamo utillizzato il generatore DC Agilent E3631A.}
\label{fig:gen_continua}
\end{wrapfigure}

Utilizzeremo ora le \textit{regole d'oro} per risolvere il circuito. B si trova a potenziale di comune, quindi per (\ref{eq:regola_V}) anche A sarà allo stesso potenziale, che considereremo nullo. Dunque è verificato
\begin{equation}
V_{gen}=I R_1
\label{eq:gen_1}
\end{equation}
Per la legge di Kirkhhoff sui nodi e dalla eq. (\ref{eq:regola_I}), si ottiene che la corrente passante per la resistenza di carico è uguale alla corrente $I$, definita dall'eq. (\ref{eq:gen_1}).

Come possiamo osservare, l'intensità di corrente passante per la resistenza $R_2$ è indipendente dal valore della resistenza stessa. Infatti la tensione di output dell'opamp verrà da esso stesso modificata in modo da far passare sempre lo stesso valore di corrente attraverso il carico.
Ciò avviene grazie al fenomeno di \textit{negative feedback}, o retroazione negativa, che permette all'opamp di controllare la tensione di output in modo tale da verificare l'eq. (\ref{eq:regola_V}).
Conoscendo la corrente attraversante $R_2$, possiamo inoltre ricavare il valore di tensione in uscita:

$$V_{out}=-\frac{R_2}{R_1} V_{gen}$$

Per quanto riguarda le misure da noi ricavate, abbiamo deciso di misurare la corrente $I$ passante per la resistenza piuttosto che la tensione di uscita, ponendo un amperometro fra l'uscita dell'opamp e la resistenza di carico $R_2$.
Abbiamo scelto di progettare il circuito affinché fornisse \SI{1}{\milli\ampere}, in modo da rimanere nel range di funzionamento ideale dell'opamp (\SIrange{10}{20}{\milli\ampere}).
Avendo a disposizione una resistenza $R_1=3.85 \pm 0.01\,\si{\kilo\ohm}$, per (\ref{eq:gen_1}), abbiamo utilizzato una tensione continua di $3.85 V$. A seguire sono riportati i valori sperimentali da noi ricavati (gli errori sono sottointesi essere unitari sull'ultima cifra significativa).

\begin{table}[H]
\center
\caption{Dati da noi registrati con relativi errori.}
{\renewcommand{\arraystretch}{1.4}%
\begin{tabular}{c|c|c|c|c|c|c|c|c}
$R_2$ [\si{\ohm}] & 0.54 & 35.1 & 412 & 1021 & 1996 & 3068 & 4170 & 4719 \\ 
\hline 
$I$ [\si{\milli\ampere}] & 1.002 & 1.002 & 1.002 & 1.002 & 1.002 & 1.002 & 1.002 & 1.002 \\ 
\end{tabular}}
\end{table}


È evidente che il circuito da noi creato fornisce un valore di corrente indipendente dal carico ad esso collegato, quantomeno nel range da noi testato.

\subsection{Sommatore Pesato}

Il secondo circuito che abbiamo preso in esame è un sommatore pesato: nel nostro caso il circuito aquisiva due segnali in ingresso e produceva in uscita un unico segnale ottenuto dalla somma pesata dei due ingressi. In particolare si osserva che il peso di ciascun segnale è dato dal rapporto tra la resistenza di feedback $R_f$ e la resistenza associata al canale di ingresso, $R_1$ e $R_2$. In Figura \ref{fig:sommatore_pesato} è riportato lo schema del circuito.

\begin{wrapfigure}[21]{l}{0.55\textwidth}
  \begin{center}
    \includegraphics[width=0.40\textwidth]{../E01/latex/c2.pdf}
  \end{center}
  \caption{Schema del sommatore pesato. Come valori abbiamo utilizzato $R_f \approx R_1=(99.9 \pm 0.1)\,\si{\kilo\ohm}$ e $R_2=(49.8 \pm 0.1)\,\si{\kilo\ohm}$, dove per $R_2$ è stato necessario utilizzare un parallelo di due resistenze da \SI{100}{\kilo\ohm}. Come GEN 1 abbiamo utilizzato l'oscilloscopio, mentre per GEN 2 il generatore di forme d'onda. Infine, per valutare la tensione in uscita abbiamo utilizzato l'oscilloscopio.}
\label{fig:sommatore_pesato}
\end{wrapfigure}

Per risolvere il circuito utilizziamo ancora una volta le \textit{regole d'oro} e le leggi di Kirchhoff. Definendo le tensioni dei generatori 1 e 2 relativamente $V_1$ e $V_2$, si ricavano le seguenti equazioni:
$$V_1 - V_A =I_1 R_1 \qquad V_2 - V_A =I_2 R_2$$
$$V_A - V_{out} =(I_1+I_2) R_f$$
Essendo però B connesso al potenziale di comune, da (\ref{eq:regola_V}) si ottiene che $V_A = V_B = 0$ che sostituita nelle equazioni precedenti porta a:
\begin{equation}
V_{out}=-R_f \left( \frac{V_1}{R_1}+\frac{V_2}{R_2}\right)
\label{eq:sum_v_out}
\end{equation}

Analizzando il caso generale di sommatore a cui sono connessi \textit{n} canali in ingresso si può riscrivere la formula (\ref{eq:sum_v_out}) più elegantemente. Definendo il peso relativo di ogni segnale come $\phi_i = \frac{R_f}{R_i}$:
\begin{equation*}
V_{out}=-\sum^{n}_{i=1} \frac{R_f}{R_{i}}V_{i}=-\sum^{n}_{i=1} \phi_i V_{i}
\end{equation*}

Noi abbiamo optato per valori semplici dei pesi: $\phi_1=1$ e $\phi_2=2$. Per ottenere questi valori abbiamo utilizzato i seguenti valori di resistenza: $R_f=R_1=\SI{100}{\kilo\ohm}$ e $R_2=\SI{50}{\kilo\ohm}$.
Presentiamo ora i grafici di alcune forme d'onda in uscita.

\begin{figure}[ht]
 \centering
   {\includegraphics[width=16.5cm]{../E01/latex/sinsin.pdf}}
 \caption{Grafico della tensione in uscita. Il generatore 1 (generatore dell'oscilloscopio) produce un'onda sinusoidale di $\nu=800$ \si{\hertz} e $V^1_{pp}=500$ \si{\milli\volt}; il generatore 2 (generatore di forme d'onda) produce un'onda sinusoidale di $\nu=100$ \si{\hertz} e $V^2_{pp}=1000$ \si{\milli\volt}. Si può verificare facilmente che l'ampiezza massima è pari a $\phi_1 V^1_{pp}+\phi_2 V^2_{pp}=2500$ \si{\milli\volt}.}
 \label{gr:onde1}
\end{figure}

\begin{figure}[ht]
 \centering
   {\includegraphics[width=16.5cm]{../E01/latex/sinquad.pdf}}
 \caption{Grafico della tensione di uscita. Il generatore 1 (generatore dell'oscilloscopio) produce un'onda sinusoidale di $\nu=900$ \si{\hertz} e $V^1_{pp}=500$ \si{\milli\volt}; il generatore 2 (generatore di forme d'onda) produce un'onda quadra di $\nu=100$ \si{\hertz} e $V^2_{pp}=1000$ \si{\milli\volt}. Anche in questo caso si può verificare facilmente che l'ampiezza massima è pari a $\phi_1 V^1_{pp}+\phi_2 V^2_{pp}=2500$ \si{\milli\volt}.}
 \label{gr:onde2}
\end{figure}
\clearpage
%\section{24.09.2014 - Amplificatori reali}

Scopo di questa esperienza è quello di studiare un amplificatore operazionale $\mu a 741$ reale.
Ne analizzeremo l'offset e le correnti di polarizzazione (\textit{bias currents}) cercando di stimarne un valore, tramite circuiti progettati ad hoc.
Premettiamo che il circuito di alimentazione è lo stesso utilizzato nella precedente esperienza e dunque non ripeteremo le considerazioni e gli schemi circuitali già proposti.
Inoltre ricordiamo che la circuiteria di alimentazione sugli schemi è stata nascosta per facilitarne la comprensione.

\subsection{Strumenti e materiali}

\begin{itemize} [noitemsep]
%\item Oscilloscopio Agilent DSO-X 2002A (bandwidth \SI{70}{\mega\hertz}, sample rate \num{2} GSa/s);
\item Generatore di tensione continua Agilent E3631A (max $\pm \, \SI{25}{\volt}$ o $\pm \, \SI{6}{\volt}$);
%\item Generatore di forme d'onta Agilent 33120A con range di frequenza da \SI{100}{\micro\hertz} a \SI{15}{\mega\hertz};
\item Multimetro Agilent 34410A a sei cifre e mezza;
\item Un amplificatore operazionale $\mu$A741;
\item Resistenze e capacità di vari valori;
\item un trimmer a un giro da \SI{5}{\kilo\ohm} e uno da \SI{10}{\kilo\ohm};
\item un trimmer multigiro da \SI{10}{\kilo\ohm};
\item Breadboard e cablaggi vari.
\end{itemize}

\subsection{Conclusioni}
In questa esperienza abbiamo potuto osservare come gli OPAMP, sebbene siano dei circuiti abbastanza precisi, abbiano delle imperfezioni, date dalla loro composizione circuitale (sono presenti dei transistor BJT al loro interno).
Le discrepanze tra il modello ideale e l'OPAMP reale sono date principalmente dallo sbilanciamento della risposta dello stesso ($V_{offset}$) e dalle correnti di polarizzazione (\textit{bias currents}).
Per nostra fortuna spesso gli OPAMP presentano dei connettori predisposti a minimizzare la tensione di offset con circuiti di compensazione: nel nostro caso un trigger collegato ai piedini di offset e all'alimentazione negativa.
Una volta bilanciato l'opamp, abbiamo misurato le correnti di polarizzazione e abbiamo potuto osservare che esse sono dell'ordine dei \si{\nano\ampere}, quindi trascurabili per gli utilizzi più comuni.

\section{Amplificatori Operazionali Reali - 24.09.2014}

In questa esperienza analizzeremo le principali caratteristiche che distinguono un amplificatore operazionale \si{\micro}A741 reale dal suo modello ideale.
In particolare andremo a studiare la tensione di offset e le correnti di polarizzazione - o \textit{bias currents} - cercando di stimare un valore per entrambe, tramite circuiti progettati ad hoc.

Prima di tutto, osserviamo che la circuiteria di alimentazione dell'opamp è la medesima della precedente esperienza e pertanto essa sarà nascosta negli schemi per facilitarne la comprensione.

\subsection{Strumenti e materiali}

\begin{itemize} [noitemsep]
\item Generatore di tensione continua Agilent E3631A (max $\pm \, \SI{25}{\volt}$ o $\pm \, \SI{6}{\volt}$);
\item Multimetro Agilent 34410A a sei cifre e mezza;
\item Un amplificatore operazionale $\mu$A741;
\item Resistenze e capacità di vari valori;
\item un trimmer a un giro da \SI{5}{\kilo\ohm} e uno da \SI{10}{\kilo\ohm};
\item un trimmer multigiro da \SI{10}{\kilo\ohm};
\item Breadboard e cablaggi vari.
\end{itemize}

\subsection{Stima e correzione dell'offset}
\label{par2:offset}

Un opamp ideale possiede una risposta in uscita perfettamente simmetrica rispetto al caso in cui i due terminali invertente e non invertente sono allo stesso potenziale.
In un opamp reale però, a causa dell'impossibilità di costruire transistor BJT con caratteristiche identiche in modo tale che la circuteria interna dell'integrato risulti simmetrica, è verificata l'esistenza di uno sbilanciamento della risposta in uscita.
La tensione di \textit{offset} $V_{os}$ è definita come la differenza di potenziale tra i terminali di entrata necessaria affinché l'uscita dell'opamp sia nulla.

\subsubsection{Configurazione senza retroazione}

La cosa più semplice da fare per osservare questa imperfezione dell'integrato è collegare entrambi i terminali di ingresso a comune (vedi Figura \ref{cir:open_loop}) e osservare grazie al multimetro la tensione in uscita.
Essa sarà data dalla tensione di offset, amplificata dal guadagno a maglia aperta $A_{ol}$ dell'opamp.
Poiché $A_{ol}$ è un valore particolarmente alto, sebbene $V_{os}$ sia relativamente piccola, spesso accade che l'opamp vada in saturazione.

Nel nostro caso l'opamp andava in saturazione negativa (\SI{-12.9}{\volt}).
Da questo dato si possono trarre due considerazioni: la prima è che il circuito si comporta come se la la tensione all'ingresso invertente fosse maggiore di quella all'ingresso non invertente.
La configurazione a maglia aperta non permette però una stima efficace del valore di $V_{os}$.

La seconda osservazione consiste nel fatto che il modulo della tensione in uscita è inferiore a quello dell'alimentazione usata (\SI{\pm15}{\volt}).
Ciò può essere dovuto alla caratteristica dell'integrato di avere il valore di tensione massimo $|V_{out,max}|$ leggermente inferiore a quello di alimentazione oppure al fatto che $V_{os}$ non è sufficiente per far andare il circuito completamente in saturazione.
Altri gruppi di laboratorio hanno infatti ricavato tensioni di molto inferiori alla tensione di alimentazione del circuito ($\lesssim \SI{10}{\volt}$).

\begin{figure}[ht]
        \centering
        \begin{subfigure}[b]{0.35\textwidth}
                 \includegraphics[width=0.70\textwidth]{../E02/latex/open_loop.pdf}
                \caption{Circuito a maglia aperta}
                \label{cir:open_loop}
        \end{subfigure}%
    \quad
        \begin{subfigure}[b]{0.35\textwidth}
               \includegraphics[width=0.70\textwidth]{../E02/latex/inv.pdf}
                \caption{Circuito amplificatore}
                \label{cir:inv}
        \end{subfigure}
     
\end{figure}

Con il circuito in Figura (\ref{cir:open_loop}) ci è pertanto impossibile ricavare una stima del valore di offset.
Dobbiamo perciò ricorrere a un circuito amplificatore che ci permetta di avere un controllo preciso sul guadagno.

\subsubsection{Configurazione con retroazione}

Lo schema in Figura (\ref{cir:inv}) rappresenta il circuito da noi usato per stimare il valore di $V_{os}$.

Trattiamo per primo il caso del circuito di amplificazione invertente.
L'idea è quella di assumere la presenza di un generatore di tensione che crei una d.d.p. $V_{os}$ tra il terminale invertente e il punto A (l'orientamento del generatore è ininfluente).
In tal modo abbiamo $V_B=0$ e quindi $V^-=V_B=0$, pertanto $V_A = V_{os}$.

Vale allora che, uguagliando le correnti nel nodo A

$$\frac{V_{os}}{R_1} + \frac{V_{os}-V_{out}}{R_2} = 0$$

da cui si ottiene
$$V_{out}=\left(1+\frac{R_2}{R_1}\right) V_{off}$$

Analogamente si tratta il caso non invertente, cioè si assume l'esistenza di un generatore di tensione che crei una d.d.p. $V_{os}$ tra comune (B) e il terminale non invertente.
Anche in questo caso l'orientamento è ininfluente.

Si avrà pertanto $V_B=0$, $V^+=V_{os}$ e pertanto $V_A=V^+=V_{os}$.
L'analisi risulta dunque identica a quella del circuito un amplificatore non invertente, e si ottiene lo stesso risultato del caso invertente, a meno di un fattore $-1$ derivante appunto dall'orientamento del generatore.

Riportiamo nella seguente tabella i valori di offset calcolati, definendo il guadagno come Gain$=1+R_2/R_1$. 

\begin{center}
\begin{savenotes}
\begin{tabular}{c|c|c|c|c}
$R_1$ [\si{\ohm}] & $R_2$ [\si{\kilo\ohm}] & Gain &$V_{out}$ [\si{\milli\volt}] & $V_{off}$ [\si{\milli\volt}]\\ 
\hline 
$119.8\pm0.1$ & $9.911\pm0.001$ & $83.73\pm0.07$&  $-103.5 \pm 0.5$ & $-1.23 \pm0.01$\\
\hline
$119.8\pm0.1$ & $99.35\pm0.01$ & $830.3\pm0.7$ &$ -1025 \pm 2$ & $-1.2 \pm0.1$\\

\end{tabular}
\end{savenotes}
\end{center}

\begin{wrapfigure}[17]{l}{0.55\textwidth}
  \begin{center}
    \includegraphics[width=0.280\textwidth]{../E02/latex/trimmer_correction.pdf}
  \end{center}
  \caption{Circuito a guadagno unitario. Il trimmer è collegato ai piedini 1 e 5 dell'opamp per compensare l'offset.}
  \label{cir2:trimmer}
\end{wrapfigure}

%Come sappiamo la tensione di offset non è però l'unico problema che incontriamo quando usiamo opamp reali. Infatti ingresso invertente e non invertente sono collegati alle basi di transistor e, ovviamente, per polarizzarli serve una corrente di base. Gli effetti di tale corrente si sommeranno dunque a quelli dovuti all'offset. Tale argomento sarà comunque trattato approfonditamente nella sezione successiva.
Per risolvere il problema della tensione di offset, si può agire collegando una resistenza variabile ai piedini 1 e 5 dell'integrato e all'alimentazione negativa $V^- = \SI{-15}{\volt}$, come mostrato in Figura \ref{cir2:trimmer}.
L'integrato è stato progettato in modo tale che, regolando questa resistenza variabile data nel nostro caso da un \textit{trimmer resistivo}, si vada a generare una tensione che bilanci l'offset.

Durante l'esperienza abbiamo provato ad utilizzare trimmer ad un giro da \SI{10}{\kilo\ohm} ma la sensibilità meccanica era troppo bassa per poter azzerare l'offset (la tensione di uscita infatti passava da $\approx$\SI{-13}{\volt} a $\approx$\SI{+13}{\volt}) troppo velocemente.
Abbiamo dunque utilizzato un trimmer multigiro da \SI{10}{\kilo\ohm}, con il quale abbiamo raggiunto la tensione $V_{out}= (1.3\pm0.2)\si{\volt}$.
Ricordando che il guadagno a maglia aperta è di 100-120dB, si nota il buon bilanciamento della tensione di offset.

\begin{wrapfigure}[17]{r}{0.55\textwidth}
  \begin{center}
    \includegraphics[width=0.280\textwidth]{../E02/latex/current_correction.pdf}
  \end{center}
  \caption{Circuito utilizzato per la stima di $V_{off}$ con la correzione sulle correnti data da $R_C$.}
  \label{cir2:current_correction}
\end{wrapfigure}

Come già accennato le correnti di bias, per quanto piccole (\si{\nano\ampere}) giocano comunque un ruolo sull'offset totale. Per minimizzare il loro impatto sul valore della tensione di offset abbiamo inserito fra l'ingresso non invertente e ground una resistenza di compensazione $R_C$ tale da annullare il contributo delle correnti di offset sulla tensione di uscita.

Il contributo alla tensione di uscita data dalle correnti di polarizzazione è

\begin{equation}
V_{out}^{C} = \left( 1+\frac{R_2}{R_1} \right)\left[ \frac{I_{b^-}R_2}{\frac{R_1+R_2}{R_1}} - I_{b^+} R_B\right]
\label{eq2:Vout_currents}
\end{equation}

dalla quale si ricava che, supponendo $I_{off} = |I_{b^+}|-|I_{b^-}| \approx 0$ e ponendo $V_{out}^{C}=0$,

$$R_C=\frac{1}{\frac{1}{R_1} + \frac{1}{R_2}} $$

cioè il parallelo fra $R_1$ ed $R_2$.

Riportiamo nella seguente tabella i nuovi valori in questa nuova configurazione (circuito in Figura \ref{cir2:current_correction}):

\begin{center}
\begin{tabular}{c|c|c|c|c|c|c}
$R_C [\si{\ohm}]$& $R_1[\si{\ohm}]$ & $R_2[\si{\kilo\ohm}]$ & Gain & $V_{out}' [\si{\milli\volt}]$ & $V_{off}' [\si{\milli\volt}]$ & $|V_{off}-V_{off}'|[\si{\milli\volt}]$ \\ 
\hline 
$119.4\pm0.1$ & $119.8\pm0.1$ & $9.911\pm0.001$  & $83.73 \pm 0.07$ & $-105.5 \pm 0.5$ & $-1.26 \pm0.01$ & $0.02\pm0.01$ \\
\hline
$119.4\pm0.1$ & $119.8\pm0.1$ & $99.35\pm0.01$  & $830.3\pm0.7$ &$ -1038 \pm 5$ & $-1.2 \pm 0.1$ & $\approx 0$\\
\end{tabular}
\end{center}

Notiamo che le differenze fra i valori della tensione di offset con e senza $R_C$ sono compatibili con il rumore ambientale di fondo (qualche decina di \si{\micro\volt} di ampiezza).

\subsection{Correnti di polarizzazione}

In questa parte dell'esperienza abbiamo progettato diversi circuiti per misurare la corrente di polarizzazione per entrambi gli ingressi, avendo già stabilizzato la tensione di offset dall'esterno. Di seguito proponiamo due modalità.

\subsubsection{Configurazione senza retroazione}

\begin{wrapfigure}[16]{r}{0.55\textwidth}
  \begin{center}
    \includegraphics[width=0.30\textwidth]{../E02/latex/direct_measure.pdf}
  \end{center}
  \caption{Schema del circuito non retro-azionato utilizzato per stimare la corrente di polarizzazione. La resistenza utilizzata è $R_B=10.36\pm0.01$\si{\mega\ohm}; la capacità $C_B=102 \pm 1$ \si{\nano\farad}.}
  \label{circuito:rel2_correnti_senzaretroazione}
\end{wrapfigure}

Nel circuito mostrato in figura abbiamo posto la resistenza all'ingresso non invertente (analogamente si può fare con l'ingresso invertente) e, misurando la caduta di potenziale ai capi della stessa con il multimetro, possiamo ottenere il valore di corrente desiderato applicando semplicemente la legge di Ohm

$$V=I_{b^+} R$$

Per far ciò, dato che attendevamo una corrente dell'ordine dei \si{\nano\ampere}, abbiamo utilizzato una resistenza molto grande in modo da poter leggere il valore della tensione su una scala accettabile per il multimetro.

Durante la procedura abbiamo però notato che, a causa di rumori ambientali, il valore di tensione sul multimetro fluttuava sulla prima cifra, rendendo nostra misurazione ovviamente non quantitativa (al massimo poteva stimarci l'ordine di grandezza della corrente). Per ovviare, abbiamo inserito in parallelo alla resistenza un condensatore che caricandosi si portava alla stessa ddp dei capi della resistenza. In questo modo abbiamo potuto ottenere un valore meno fluttuante, che si attestava a $V=(-80 \pm 2)$ \si{\milli\volt}, cioè $I_{b^+}=(7.7 \pm 0.2)$ \si{\nano\ampere}.

Con questo metodo semplice abbiamo potuto ottenere una prima stima del valore della corrente. Di contro bisogno considerare che il rumore non permette di avere una stima qualitativa ed inoltre la resistenza, scaldandosi, modifica il suo valore e potrebbe portare ad un errore sulla misura. Successivamente progetteremo dunque un circuito che, sfruttando l'amplificazione data dall'amplificatore operazionale, minimizzerà questi errori.

Analizziamo ora una possibile causa della  fluttuazione della misura.

\subsubsection*{Rumore termico}

Abbiamo considerato che la fluttuazione della misura sia dovuta ad un possibile effetto del rumore termico sulla nostra resistenza da \SI{10}{\mega\ohm}. Considerando la larghezza di banda del multimetro $\Delta f = \SI{1}{\hertz}$ e siano $k_B = \SI{1.38e-23}{\joule\per\kelvin}$ la costante di Boltzmann, $T$ la temperatura assoluta e $R$ il valore della resistenza presa in considerazione

\begin{equation}
	V_{eff}^2 = 4 k_B T R \Delta f
\end{equation}

da cui, ovviamente, si ottiene che il possibile rumore termico è 5 ordini di grandezza inferiore della tensione misurata

\begin{equation}
	V_{eff} = \sqrt{4 k_B T R \Delta f} \simeq 4 \times 10^{-7} \si{\volt}
\end{equation}

ed è pertanto ininfluente. Dunque il rumore presente è sicuramente imputabile ad altre cause ambientali.

\subsubsection{Configurazione con retroazione negativa}

Sfruttando un modello simile a quello utilizzato per trovare la tensione di offset, abbiamo montato i circuiti come in figura. Data la tensione in uscita, grazie alle proprietà di amplificazione dei segnali in ingresso dell'OPAMP, possiamo ottenere una misura indiretta della corrente di polarizzazione.

\subsubsection*{Misura di $I_{b^-}$}

\begin{wrapfigure}[15]{l}{0.55\textwidth}
  \begin{center}
    \includegraphics[width=0.30\textwidth]{../E02/latex/inv_current.pdf}
  \end{center}
  \caption{Schema del circuito retro-azionato utilizzato per stimare la corrente di polarizzazione $I_{b^-}$. Le resistenze utilizzate sono $R_1=(98.9\pm0.1)$ \si{\ohm}, $R_2=(99.4\pm0.1)$ \si{\kilo\ohm} e $R_B=(99.4\pm0.1)$ \si{\kilo\ohm}.}
  \label{circuito:rel2_correnti_retroazione_inv}
\end{wrapfigure}

Risolviamo il circuito per trovare la corrente di polarizzazione $I_{b^-}$ in funzione della tensione di uscita. Considerando $V_{-}$ la tensione al capo di $R_B$ collegato all'OPAMP e $V^*$ quello opposto, vale in quel punto la legge di Kirchhoff sui nodi

$$\frac{V^* - V_{in}}{R_1} + \frac{V^*-V_{out}}{R_2} + \frac{V^*-V_{-}}{R_B}=0$$

Dato che l'amplificatore operazionale è considerato già stabilizzato per quanto riguarda la tensione di offset, possiamo considerare la tensione all'ingresso invertente uguale all'ingresso non invertente. Vale dunque che $V_{in}=V_{-}=0$ e si trova (considerando $I_{b^-} R_B = V^*$):

$$I_{b^-}=\frac{V_{out}}{R_2 R_B}\frac{1}{\frac{1}{R_1}+\frac{1}{R_2}+\frac{1}{R_B}}$$

Le resistenze sono state dimensionate tenendo invece conto dell'ordine di grandezza della corrente da misurare e considerando il risultato sopra ottenuto: volevamo che $V_{out}$ fosse almeno $10^7\approx10^8$ volte più grande della corrente, per poter utilizzare il multimetro, che ha scale di misura limitate. I valori sono in Figura \ref{circuito:rel2_correnti_retroazione_inv}.

La misura di tensione di uscita è di $(3.89\pm0.02)$ \si{\volt} ed il valore ottenuto è dunque $I_{b^-} = (38 \pm 5)$ \si{\nano\ampere}.

\subsubsection*{Misura di $I_{b^+}$}

\begin{wrapfigure}[17]{r}{0.55\textwidth}
  \begin{center}
    \includegraphics[width=0.25\textwidth]{../E02/latex/ninv_current.pdf}
  \end{center}
  \caption{Schema del circuito retro-azionato utilizzato per stimare la corrente di polarizzazione $I_{b^+}$. Le resistenze utilizzate sono le medesime del circuito precedente in Figura \ref{circuito:rel2_correnti_retroazione_inv}.}
  \label{circuito:rel2_correnti_retroazione_noninv}
\end{wrapfigure}

Similmente a quanto visto per la configurazione prima, troviamo che, data la legge di Kirchhoff (con $V^*$ la tensione all'ingresso non invertente, che per quanto detto sopra è uguale a quella all'ingresso invertente)

$$\frac{V^* - V_{in}}{R_1} + \frac{V^*-V_{out}}{R_2}=0$$

e considerando $V^*=I_{b^+} R_B$, otteniamo

\begin{equation}
I_{b^+}=\frac{V_{out}}{R_2 R_B}\frac{1}{\frac{1}{R_1}+\frac{1}{R_2}}
\label{eq2:corrente_noninv}
\end{equation}

Anche in questo caso le resistenze sono state dimensionate come sopra e i valori sono in Figura \ref{circuito:rel2_correnti_retroazione_noninv}. La misura di tensione di uscita è di $-(3.72 \pm 0.02)$   \si{\volt} ed il valore ottenuto è dunque $I_{b^+} = - (37.2 \pm 0.2)$ \si{\nano\ampere} \footnote{La negatività della corrente va intesa rispetto all'ingresso non invertente, ed è quindi uscente rispetto a tale ingresso. Al contrario, nel paragrafo precedente, la corrente è intesa entrante nel punto di $V^*$, e quindi è entrante rispetto all'ingresso invertente.}.

\subsubsection*{Calcolo di $I_{b^+}$ data $I_{b^-}$}

Consideriamo ora un altro modo per trovare la corrente di polarizzazione $I_{b^+}$ supponendo di aver già effettuato la misura di $I_{b^-}$ nel primo circuito (Figura \ref{circuito:rel2_correnti_retroazione_inv}). Successivamente controlleremo che il valore 'sperimentale' calcolato con (\ref{eq2:corrente_noninv}) è compatibile con quello calcolato in questo paragrafo.

Analizziamo dunque il secondo circuito (Figura \ref{circuito:rel2_correnti_retroazione_noninv}) per cercare di trovare la dipendenza di $I_{b^+}$ da $I_{b^-}$. Vale, dalla teoria, l'equazione (\ref{eq2:Vout_currents}), da cui è possibile ricavare

$$I_{b^+} = \frac{R_1}{R_B(R_1+R_2)}(I_{b^-} R_2-V_{out})$$

ed inserendo i valori otteniamo $I_{b^+} = - (37.2 \pm 0.2)$ \si{\nano\ampere}, compatibile con il risultato precedente.

\subsection{Conclusioni}
In questa esperienza abbiamo potuto osservare come gli OPAMP, sebbene siano dei circuiti abbastanza precisi, abbiano delle imperfezioni, date dalla loro composizione circuitale (sono presenti dei transistor BJT al loro interno).
Le discrepanze tra il modello ideale e l'OPAMP reale sono date principalmente dallo sbilanciamento della risposta dello stesso ($V_{offset}$) e dalle correnti di polarizzazione (\textit{bias currents}).
Per nostra fortuna spesso gli OPAMP presentano dei connettori predisposti a minimizzare la tensione di offset con circuiti di compensazione: nel nostro caso un trigger collegato ai piedini di offset e all'alimentazione negativa.
Una volta bilanciato l'opamp, abbiamo misurato le correnti di polarizzazione e abbiamo potuto osservare che esse sono dell'ordine dei \si{\nano\ampere}, quindi trascurabili per gli utilizzi più comuni.
\clearpage
\section{Amplificatori Operazionali Reali - Seconda Parte - 30.09.2014}

In questa esperienza abbiamo studiato ulteriori caratteristiche che distinguono l'opamp reale dal modello ideale.
Più precisamente abbiamo analizzato e misurato lo \textit{slew rate}, la massima corrente pilotabile, la larghezza di banda passante e il guadagno a maglia aperta $A_{ol}$ di un opamp $\mu$A741.

Ricordiamo che in ogni schema sono stati sottintesi la circuitazione di alimentazione (cavi e capacitori) e quella dedicata al bilanciamento dell'offset (cavi e trimmer).

\subsection*{Materiali}

\begin{itemize} [noitemsep]
%\item Oscilloscopio Agilent DSO-X 2002A (bandwidth \SI{70}{\mega\hertz}, sample rate \num{2} GSa/s);
%\item Generatore di tensione continua Agilent E3631A (max $\pm \, \SI{25}{\volt}$ o $\pm \, \SI{6}{\volt}$);
%\item Generatore di forme d'onta Agilent 33120A con range di frequenza da \SI{100}{\micro\hertz} a \SI{15}{\mega\hertz};
%\item Multimetro Agilent 34410A;
\item Un amplificatore operazionale $\mu$A741;
\item Resistenze e capacità di vari valori;
%\item un trimmer (potenziometro);
%\item Breadboard e cablaggi vari.
\end{itemize}

\subsection{Misura dello Slew Rate}

\begin{wrapfigure}[15]{r}{0.55\textwidth}
  \begin{center}
    \includegraphics[width=0.32\textwidth]{../E03/latex/slew_rate.pdf}
  \end{center}
  \caption{Schema del circuito utilizzato per stimare lo Slew Rate. La resistenza utilizzata è $R=\SI{2.17\pm0.01}{\kilo\ohm}$; la capacità $C=\SI{200 \pm 10}{\pF}$.}
  \label{cir3:slew_rate}
\end{wrapfigure}

Lo Slew Rate è uno dei parametri che caratterizzano l'amplificatore operazionale ed è definito come
$$SR = \mathrm{max}\left[\frac{\Delta V}{\Delta t}\right]$$
Esso indica la massima velocità con cui l'opamp varia la tensione in uscita nell'unità di tempo.
Ciò significa che, per segnali in entrata con derivata temporale superiore ad SR dell'amplificatore operazionale, avremo un segnale distorto in uscita.

Per valutare la risposta in uscita del nostro $\mu$A741 abbiamo utilizzato il circuito in Figura \ref{cir3:slew_rate} fornendo in ingresso un'onda quadra caratterizzata da una velocità di variazione dei segnale di $SR_{gen}=\SI{315}{\volt\per\micro\second}$ (misurata tramite l'oscilloscopio).
Essendo $SR_{gen}$ diversi ordini di grandezza maggiore del valore atteso di $SR$ ($\approx \SI{0.5}{\volt\per\micro\second} $) possiamo considerare l'onda quadra come ideale, cioè avente pendenza infinita.
La capacità posta in parallelo alla resistenza è utilizzata come serbatoio di cariche per attenuare fenomeni di rumore nel segnale creati dal generatore.

\begin{figure}[ht]
 \centering
   {\includegraphics[width=\textwidth]{../E03/latex/sr_uad.pdf}}
 \caption{Grafico della tensione in uscita (rosso) ed in entrata (nero) in funzione del tempo. I $\Delta$ sono relativi alle differenze fra le coordinate dei punti segnati. Si noti che le scale dei tempi sono differenti per i due grafici.}
 \label{gr3:slew_rate}
\end{figure}

Abbiamo posto inizialmente la frequenza a $f=1$\si{\kilo\hertz} e la tensione picco-picco a $V_{pp}=10$ \si{\volt};
abbiamo poi misurato con i cursori dell'oscilloscopio la differenza fra il 10\% e il 90\% del valore di tensione picco picco in uscita $\Delta V = (8.07 \pm 0.01)$ \si{\volt}  e la relativa $\Delta t = (15.9 \pm 0.1)$ \si{\micro\second}, ottenendo uno $SR= \frac{\Delta V}{\Delta t} = \SI{0.507 \pm 0.003}{\volt\per\micro\second}$.

Però, come si può notare dal grafico in Figura \ref{gr3:slew_rate}, nella zona vicina al 10\% è presente una parte di segnale affetta da rumore che potrebbe rendere poco precisa la nostra misura.
Per ovviare a questo problema abbiamo spostato il limite inferiore ad un valore percentuale più alto (25\%).
Quest'operazione è permessa dal fatto che siamo interessati al rapporto fra i $\Delta$ e questo rimane costante su una retta.
Portando dunque la percentuale del limite inferiore al 25\% abbiamo ottenuto i seguenti valori: $\Delta V = \SI{6.48 \pm 0.01}{\volt}$ e $\Delta t = \SI{14.5 \pm 0.1}{\ms}$, dunque $SR = \SI{0.447 \pm 0.003}{\volt\per\ms}$. Durante l'esperienza utilizzeremo questo valore come riferimento per lo slew rate del nostro amplificatore, in quanto più restrittivo.

Abbiamo anche notato che lo slew rate è differente a seconda che il segnale abbia derivata $dV/dt$ positiva o negativa (grafico in Figura \ref{gr3:slew_rate}, immagine di destra).
Nel secondo caso, infatti, abbiamo ottenuto i seguenti valori: $\Delta V = (8.00 \pm 0.01)$ \si{\volt} e $\Delta t = (23.2 \pm 0.1)$ \si{\micro\second}, quindi $SR = (0.345 \pm 0.004)$ \si{\volt\per\micro\second}.

\subsection{Misura della capacità di pilotaggio}

\begin{wrapfigure}[18]{l}{0.55\textwidth}
  \begin{center}
    \includegraphics[width=0.33\textwidth]{../E03/latex/max_current.pdf}
  \end{center}
  \caption{Schema del circuito utilizzato per misurare la corrente massima. La resistenza utilizzata è $R=\SI{98.47\pm0.01}{\ohm}$; il generatore fornisce un'onda triangolare di $f=1$ \si{\kilo\hertz} e $5 V_{pp}$.}
  \label{cir3:max_current}
\end{wrapfigure}

Un altro limite dell'amplificatore operazionale è quello di avere un valore massimo di corrente erogabile in uscita.
Pertanto, utilizzando un circuito di feedback come in Figura \ref{cir3:max_current}, la corrente massima e il valore della resistenza $R$ fisseranno la tensione letta ai capi della resistenza stessa, secondo questa relazione:
\begin{equation*}
	I_{max} = \frac{V_{clip}}{R}
\end{equation*}
dove $V_{clip}$ è la tensione in cui il segnale in uscita risulta tagliato rispetto a quello in ingresso (fenomeno del \textit{clipping}).

Invertendo la precedente equazione, abbiamo dunque potuto misurare $I_{max}$, ponendo una resistenza di carico fra l'uscita dell'operazionale e terra\footnote{Si noti che, data l'alta impedenza in ingresso dell'oscilloscopio (\SI{1}{\mega\ohm}), non era possibile misurare la corrente massima con questo strumento (attesa, come da specifiche del costruttore, sui \SI{15}{\milli\ampere}).}, con l'oscilloscopio ai capi di tale resistenza per la misura di tensione (circuito in Figura \ref{cir3:max_current}).

\begin{figure}[ht]
 \centering
   {\includegraphics[width=\textwidth]{../E03/latex/clip.pdf}}
 \caption{Grafico della tensione in uscita (verde) ed in entrata (nero) in funzione del tempo. Si nota subito che il clip a tensioni negative avviene a tensioni in modulo maggiori.}
 \label{gr3:clip}
\end{figure}

Bisogna però considerare che, essendo il segnale in entrata in tensione alternata, la corrente risulta oscillare fra valori negativi e positivi (corrente rispettivamente entrante o uscente dal nodo A verso $R$.
Si può notare anche solo osservando l'asimmetria nel grafico in Figura \ref{gr3:clip} che i valori per le due correnti sono in modulo differenti.
Se consideriamo $V^+ = \SI{1.92 \pm 0.01}{\volt}$ la tensione di clip positiva e $V^- = \SI{3.45 \pm 0.01}{\volt}$ quella negativa, otteniamo che le correnti massime erogabili dall'opamp sono:

$$I_{V^+} = \frac{V^+}{R} = \SI{19.5 \pm 0.1}{\mA}  \qquad I_{V^-} = \frac{V^-}{R} = \SI{35.2 \pm 0.2}{\mA}$$
entrambi i valori sono compatibili con quelli del costruttore (valori di corrente fino a $\pm 40$ \si{\milli\volt}).

L'asimmetria delle correnti può essere spiegata considerando l'asimmetria della circuiteria interna dell'opamp, la quale è composta da diversi transistor che solo idealmente possono essere tutti delle medesime ed identiche caratteristiche.

\newpage
\subsection{Verifica della banda passante}
\label{par3:bode}

In questa parte dell'esperienza vogliamo valutare la banda passante di un circuito amplificatore (non invertente) considerando che il guadagno a maglia aperta dell'amplificatore operazionale dipende della frequenza.

\subsubsection*{Funzione di trasferimento e frequenza di taglio}

Esaminiamo dapprima il circuito, calcolandone la funzione di trasferimento. Essa è definita come:\phantom{xxxxxxxxxxxxxxxxx}
\begin{wrapfigure}[20]{l}{0.45\textwidth}
  \begin{center}
    \includegraphics[width=0.32\textwidth]{../E03/latex/bandwidth.pdf}
  \end{center}
  \caption{Circuito utilizzato per valutare la banda passante. $R_1=\SI{997\pm1}{\ohm}$ è fissa; mentre $R_2$ è stata cambiata con valori $R_2^{11\mathrm{x}}=\SI{9.91 \pm 0.01}{\kilo\ohm}$ e $R_2^{101\mathrm{x}}=\SI{99.1 \pm 0.1}{\kilo\ohm}$ a seconda del guadagno desiderato.}
  \label{cir3:banda}
\end{wrapfigure}
\begin{equation}
H(s)=\frac{V_{out}}{V_{in}}
\label{eq3:def_f_trasf}
\end{equation}
dove $s=j\omega$.
Per ogni amplificatore operazionale vale
\begin{equation}
V_{out}=A(s) (V^+-V^-)
\label{eq3:regola_opamp}
\end{equation}
con il guadagno che dipende da $s=j\omega$.
Consideriamo
$$V^+ = V_{in} \qquad V^-=V_{out} \frac{R_1}{R_1+R_2}$$
dove $V^-$ è ricavato dall'analisi delle correnti sul ramo di retroazione.
Sostituendo questi valori in (\ref{eq3:regola_opamp}) otteniamo
$$A(s) V_{in} = V_{out} + V_{out} A(s) \beta \quad con \quad \beta = \frac{R_1}{R_1+R_2} = \frac{1}{G}$$
Si noti che $G$ è il guadagno circuitale di un amplificatore non invertente.
Si può ora calcolare $H(s)$:
\begin{equation}
H(s)=\frac{V_{out}}{V_{in}}=\frac{1}{\beta}\frac{1}{1+\frac{1}{A(s) \beta}}
\label{eq3:funz_trasfe}
\end{equation}
da cui è facile notare che per $A(s) \rightarrow + \infty$ si ottiene l'approssimazione di amplificatore ideale $H(s)=G=\frac{1}{\beta}=1+\frac{R_2}{R_1}$ e l'equazione diventa indipendente dal guadagno a maglia aperta $A$.

Se schematizziamo l'opamp come un filtro passa basso, abbiamo che il guadagno a maglia aperta varia con la frequenza secondo la legge
\begin{equation}
A(j\omega)=\frac{A_{ol}}{1+j\frac{\omega}{\omega_0}}
\label{eq3:passa_basso}
\end{equation}
con $A_{ol}$ il guadagno dell'operazionale per segnali costanti ($f = 0$) e $\omega_0 \approx 8$ \si{\hertz} la prima frequenza di taglio dell'operazionale data dalla capacità di compensazione nella circuiteria interna. Sostituendo questa relazione in (\ref{eq3:funz_trasfe}), abbiamo
\begin{equation}
H(j\omega)=\frac{\frac{A_{ol}}{1+A_{ol}\beta}}{1+j \frac{\omega}{(1+A_{ol}\beta)\omega_0}} \approx \frac{G}{1+j\frac{\omega}{\omega_t}}
\label{eq3:funzione_trasferimento}
\end{equation}
da cui si nota che la nuova frequenza di taglio è data da
\begin{equation}
\omega_t=(1+A_{ol}\beta)\omega_0
\label{eq3:freq_taglio}
\end{equation}
Nel circuito in Figura \ref{cir3:banda} abbiamo impostato un guadagno $G$ retroazionato di $11$ e $101$, quindi ci aspettiamo\footnote{Questi valori sono approssimati perché cercheremo di calcolare $A_{ol}$ con le frequenze a nostra disposizione, ricavate con l'analisi della banda, piuttosto che calcolare le frequenze di taglio supponendo certo $A_{ol}$.} frequenze di taglio, ricavate con l'equazione precedente, rispettivamente sull'ordine dei $10^5$ e $10^4$.

Un secondo modo per fare una stima delle frequenze di taglio è quello di sfruttare la costante \textit{gain-bandwidth product} GBWP $= \omega_t G$.

\subsubsection*{Effetto dello Slew Rate}

Durante la misurazione bisogna però tenere conto dello slew rate.
Esso infatti può influenzare negativamente la misura distorcendo la forma d'onda in uscita.
È necessario pertanto scegliere valori di tensione del segnale in ingresso appropriati.
Supponiamo dunque di avere una forma d'onda in uscita del tipo $V_{out} \sin (2 \pi f_0 t) = V_{in} G \sin (2 \pi f_0 t)$.
Affinché non si raggiunga lo SR, deve essere verificato
$$\mathrm{max}\left[ \frac{dV_{out}}{dt} \right] = \mathrm{max} \left[ 2 \pi V_{in} G f_0 \cos (2 \pi f_0 t) \right]= 2 \pi V_{in} G f_0 < 0.35 \si{\volt\per\micro\second}$$
dove l'ultimo valore è stato scelto per tenerci sempre abbastanza distanti dal valore reale di SR. Dunque vale che
\begin{equation}
V_{in}<\frac{0.35 \si{\volt\per\micro\second}}{2 \pi G f_0} = \frac{0.35 \si{\volt\per\micro\second}}{2 \pi \mathrm{GBWP}}
\label{eq3:tensione_max}
\end{equation}
Considerando $f_0$ come la frequenza di taglio stimata prima, otteniamo che i valori di tensione da usare nei due casi è di circa $V_{in}^{11\mathrm{x}} = 100$ \si{\milli\volt}$= V_{in}^{101\mathrm{x}}$. Si nota che la regola \ref{eq3:tensione_max} è valida quando il guadagno non dipende dalla frequenza: infatti dopo la frequenza di taglio il guadagno del circuito $G$ cala e dunque i valori di tensione ammessi risultano di valore maggiore.

\subsubsection*{Grafici di Bode}
\label{par3:sub_bode}

\begin{figure}[ht]
 \centering
   {\includegraphics[width=0.75\textwidth]{../E03/latex/bode.pdf}}
 \caption{Diagrammi di Bode per il guadagno del circuito in Figura \ref{cir3:banda}. Le leggi plottate sono (\ref{eq3:fit_gain}) e (\ref{eq3:fit_fase}), con i rispettivi guadagni $G$ del circuito.}
 \label{gr3:bode}
\end{figure}

Utilizzando una forma d'onda sinusoidale abbiamo dunque misurato la tensione di uscita del circuito e la fase fra i segnali, per poi plottare i grafici di Bode (Figura \ref{gr3:bode}).

Dalla funzione di trasferimento (\ref{eq3:funzione_trasferimento}) possiamo inoltre trovare le relazioni fra il guadagno e la fase con la frequenza. Sappiamo infatti che il guadagno è dato $V_{out}/V_{in}$, ed in generale queste due tensioni possono essere considerate come complesse. Se siamo dunque interessati al guadagno dobbiamo considerare piuttosto il rapporto dei moduli, ed otteniamo la seguente relazione
\begin{equation}
\mathrm{Gain}=\frac{|V_{out}|}{|V_{in}|}=|H(j\omega)|=\sqrt{\frac{A_{ol}^2 \omega_0^2}{(\omega_0 + A_{ol} G \omega_0)^2 + \omega^2}}
\label{eq3:fit_gain}
\end{equation}
mentre, per la fase utilizziamo la solita formula della fase per un numero complesso (nel nostro caso $V_{out}/V_{in}$), ed otteniamo
\begin{equation}
\varphi=\arctan\left[{\frac{\mathrm{Im}[H(j\omega)]}{\mathrm{Re}[H(j\omega)]}}\right]= - \arctan\left[\frac{\omega}{(1+A_{ol}G)\omega_0}\right]
\label{eq3:fit_fase}
\end{equation}

Date queste leggi è possibile stimare $A_{ol}$ e $\omega_0$ dell'operazionale: rimandiamo al paragrafo \ref{par3:A_ol} per questi calcoli, dove saranno confrontati i valori ottenuti con questi dati e quelli presi al paragrafo \ref{par3:open}.

Infine, i valori di frequenza di taglio possono essere calcolati dalla legge teorica nel punto di flesso ($45^{\circ}$). Questi sono: $f_t^{101\mathrm{x}} = (10.3 \pm 0.1)$ \si{\kHz} e $f_t^{11\mathrm{x}} = (92 \pm 1)$ \si{\kHz}, in accordo con gli ordini di grandezza stimati in precedenza.

\newpage
\subsection{Guadagno Open Loop}
\label{par3:open}

\begin{wrapfigure}[17]{r}{0.45\textwidth}
  \begin{center}
    \includegraphics[width=0.280\textwidth]{../E03/latex/LF_ol.pdf}
  \end{center}
  \caption{Circuito utilizzato per misurare il guadagno $A_{ol}$. I valori di $R_4$ utilizzati sono $\SI{5.2\pm0.1}{\ohm}$,  $\SI{9.6\pm0.1}{\ohm}$, $\SI{101.3\pm0.1}{\ohm}$ e $\SI{998.2\pm0.1}{\ohm}$ mentre $R_1=\SI{98\pm1}{\kohm}$, $R_2=\SI{97\pm1}{\kohm}$ e $R_3=\SI{98\pm1}{\kohm}$.}
  \label{cir3:low_frequency}
\end{wrapfigure}

Come ultima parte dell'esperienza, abbiamo cercato di misurare il guadagno \textit{open-loop} $A_{ol}$ del nostro amplificatore reale.
Come sappiamo, tale guadagno è una funzione della frequenza e dobbiamo dunque trovare un modo per misurarlo.

L'opamp $\mu$A741 è progettato per avere un guadagno costante di circa \num{2E5} fino alla frequenza di circa \SI{8}{\Hz}.
Oltre tale frequenza abbiamo una caduta di circa \SI{8}{\decibel}/ottava (\SI{20}{\decibel}/decade) ed è pertanto impossibile effettuare misure a loop aperto in quanto piccolissime variazioni di tensione ai due ingressi causerebbero grandi effetti in uscita, mandando l'integrato in saturazione.

Per questo motivo abbiamo utilizzato il circuito riportato in Figura \ref{cir3:low_frequency}, in modo da controllare la tensione in uscita e limitarla a valori inferiori alla tensione di saturazione.
Utilizzando l'oscilloscopio abbiamo poi misurato i valori di tensione $V_A$ e $V_{out}$.

Cerchiamo dunque di legare il guadagno a tali quantità.
Nel caso dell'amplificatore operazionale sappiamo che vale l'eq. (\ref{eq3:regola_opamp}), dove $(V_+-V_-)$ è semplicemente la differenza di potenziale presente tra i due ingressi, che sarà ovviamente data da $\Delta V = I_3R_4$, dove $I_3$ è la corrente che scorre attraverso $R_3$ ed $R_4$\footnote{Assumiamo che la corrente assorbita dagli ingressi sia trascurabile.}.
È quindi facile stimare $I_3=\frac{V_A}{R_3+R_4}$ conoscendo $V_A$.

Si ottiene dunque:
\begin{equation}
A_{ol}=\frac{{V_{out}}^{pp}}{{V_A}^{pp}} \frac{R_4+R_3}{R_4}
\label{eq3:lfgain}
\end{equation}

Osserviamo come in equazione (\ref{eq3:lfgain}) non compaia il termine $V_{in}$.
Dovremo dunque scegliere per le varie frequenze una tensione picco-picco in ingresso adeguata in modo che il nostro opamp non saturi e allo stesso tempo $V_A$ sia sufficientemente grande da essere poco influenzata dal rumore di fondo e sia stabile il più possibile.

Durante le misurazioni abbiamo deciso di variare la resistenza $R_4$ suddividendo il procedimento di misura in diversi step; questa strategia ci ha permesso di mantenere i valori di $V_A$ e $V_{out}$ sempre nell'ordine del volt.
I dati ricavati sono esposti nel grafico in Figura \ref{cir3:gain_open_loop}.

\subsubsection*{Misura per alte frequenze}

\begin{wrapfigure}[13]{r}{0.45\textwidth}
  \begin{center}
    \includegraphics[width=0.28\textwidth]{../E03/latex/HF_ol.pdf}
  \end{center}
  \caption{Circuito utilizzato per stimare il guadagno $A_{ol}$ a basse frequenze. I valori delle resistenze utilizzate sono le stesse del circuito in Figura \ref{cir3:low_frequency}.}
  \label{cir3:high_frequency}
\end{wrapfigure}

Per non dover utilizzare resistenze $R_4$ troppo grandi ad alte frequenze ($>$ \SI{70}{\kilo\hertz}), abbiamo provato ad effettuare le misure direttamente ad open-loop, dato che il guadagno è contenuto (\num{<100}).
Tuttavia il generatore di forme d'onda non emette un segnale pulito alla sola frequenza voluta, in quanto anch'esso è uno strumento reale e dunque affetto da errori.
Esso potrebbe dunque emettere segnali a frequenze basse o addirittura un piccolo offset.
Poiché per tali frequenze il guadagno dell'amplificatore operazionale è enorme, il segnale in uscita viene immediatamente distorto da tali rumori indesiderati.
Non è da sottovalutare poi la presenza dei residui dell'offset dell'opamp.
Esso infatti dipende anche dalla temperatura e dunque la sua compensazione non potrà mai essere perfetta ad ogni temperatura.
Misurare il guadagno $A_{ol}$ effettivamente a open-loop è particolarmente problematico.

Per aggirare tali ostacoli possiamo utilizzare il circuito riportato in Figura \ref{cir3:low_frequency}, rimuovendo la resistenza $R_4$ e sostituendo $R_3$ con un filo partendo dalla configurazione di Figura \ref{cir3:high_frequency}.
Cosi facendo otteniamo la semplice relazione $A_{ol}=\frac{{V_{out}}^{pp}}{{V_A}^{pp}}$, che possiamo anche vedere come limite di eq. (\ref{eq3:lfgain}) per $R_4 \rightarrow \infty$.
I dati ottenuti sono riportati nel grafico in Figura \ref{cir3:gain_open_loop}.

\begin{figure}[h]
	\centering
	\includegraphics[width=0.9\textwidth]{../E03/latex/gol.pdf}
	\caption{Grafico semilogaritmico del guadagno (in $dB$) in funzione della frequenza. Ad ogni forma dei punti (e relativo colore) corrisponde una diversa resistenza $R_4$ utilizzata. Per quanto riguarda $R_4=\infty$, intendiamo la mancanza del collegamento fra ingresso invertente e non invertente (circuito in Figura \ref{cir3:high_frequency}). La legge plottata è la (\ref{eq3:fit_gain_ol}).}
  \label{cir3:gain_open_loop}
\end{figure}

Anche in questo caso, considerando l'operazionale come un semplice passa basso, possiamo utilizzare l'equazione (\ref{eq3:passa_basso}) e trovare la relazione fra il guadagno e la frequenza. Otteniamo
\begin{equation}
\mathrm{Gain}=\frac{|V_{out}|}{|V_{in}|}=\sqrt{\frac{A_{ol}^2 \omega_0^2}{\omega^2 + \omega_0^2}}
\label{eq3:fit_gain_ol}
\end{equation}

Per un fit e i parametri si veda il paragrafo successivo.

\newpage
\subsection{Calcolo di $A_{ol}$}
\label{par3:A_ol}

\subsubsection{$A_{ol}$ da circuiti retroazionati}

Abbiamo ottenuto una stima di $A_{ol}$ e $\omega_0$ utilizzandoli come parametri su un fit delle leggi (\ref{eq3:fit_gain}) e (\ref{eq3:fit_fase}) sui dati sperimentali.
Inoltre, dato che l'opamp utilizzato era lo stesso per tutte le misurazioni, con una media pesata sui valori ottenuti con i quattro fit (due con (\ref{eq3:fit_gain}) e due con (\ref{eq3:fit_fase})), abbiamo ottenuto i seguenti valori sperimentali
$$A_{ol}=(1.8\pm0.1)\times 10^5 \qquad f_0=(5.9\pm0.2)\si{\hertz}$$

Ovviamente possiamo ottenere un valore di $A_{ol}$ anche utilizzando la (\ref{eq3:freq_taglio})
$$A_{ol} = \frac{G}{\omega_0} (\omega_t - \omega_0)$$
e considerando la frequenza $f_0$ ricavata dal fit. Si ottiene un valore teorico (media dei valori con i due guadagni del circuito) di $A_{ol} = (1.71\pm 0.06)\times 10^5$, compatibile con il valore sperimentale.

\subsubsection{$A_{ol}$ da open loop}

Con il fit dalla (\ref{eq3:fit_gain_ol}) otteniamo i seguenti valori
$$A_{ol}=(2.0\pm0.1)\times 10^5 \qquad f_0=(5.8\pm0.2)\si{\hertz}$$

I dati ottenuti non risultano compatibili, riteniamo che ciò sia dovuto ad errori parassiti non ben identificati.

\subsection*{Conclusioni}

In questa esperienza siamo riusciti a stimare lo "slew rate" di un amplificatore operazionale $\mu$A741 e la massima corrente erogata. Tali valori sono risultati compatibili con quelli riportati sul datasheet del costruttore. 

Successivamente abbiamo verificato la banda passante per un circuito non invertente retroazionato negativamente con guadagno G=11 e G=101. Dai diagrammi di Bode abbiamo visto come aumentando il Gain diminuiva la banda passante (ovvero GBWP=const).

Infine abbiamo valutato il guadagno open-loop a diverse frequenze, utilizzando dei circuiti che ci permettevano un controllo del guadagno effettivo e dunque di evitare problemi quali la saturazione dell'amplificatore operazionale. Il guadagno open-loop per tensioni costanti è stato stimato essere un valore compreso fra \num{1.8E5} e \num{2.0E5}, valore compatibile con quelli forniti dal costruttore.

\clearpage
\section{08.10.2014 - Comparatori}

In questa esperienza studieremo dei circuiti comparatori con retroazione positiva. In particolare, analizzeremo il funzionamento del \textit{trigger di Schmitt} non invertente, dell'oscillatore a rilassamento e di un interruttore crepuscolare. 


\subsection*{Strumenti e materiali}

\begin{itemize} [noitemsep]
\item Oscilloscopio Agilent DSO-X 2002A (bandwidth \SI{70}{\mega\hertz}, sample rate \num{2} GSa/s);
\item Generatore di tensione continua Agilent E3631A (max $\pm \, \SI{25}{\volt}$ o $\pm \, \SI{6}{\volt}$);
\item Generatore di forme d'onta Agilent 33120A con range di frequenza da \SI{100}{\micro\hertz} a \SI{15}{\mega\hertz};
\item Multimetro Agilent 34410A a sei cifre e mezza;
%\item Un amplificatore operazionale $\mu$A741;
\item Un amplificatore operazionale LM311;
\item Fototransistor OP550;
\item Resistenze e capacità di vari valori;
\item un trimmer a un giro da \SI{5}{\kilo\ohm} e uno da \SI{10}{\kilo\ohm};
\item un trimmer multigiro da \SI{10}{\kilo\ohm};
\item Breadboard e cablaggi vari.
\end{itemize}

\subsection{Verifica del funzionamento di un comparatore}

In questa prima parte dell'esperienza abbiamo verificato il funzionamento di un amplificatore operazionale LM311 come comparatore.

\begin{wrapfigure}[18]{r}{0.4\textwidth}
  \begin{center}
    \includegraphics[width=0.280\textwidth]{../E04/latex/LM311.pdf}
  \end{center}
  \caption{Piedinatura dell'operazionale LM311.}
  \label{cir4:piedinatura_LM311}
\end{wrapfigure}

Questo amplificatore ha uno slew rate molto alto e lavora sull'ordine dei nanosecondi: rispetto agli $0.5$\si{\volt\per\micro\second} del $\mu$A741, lo slew rate particolarmente elevato permette di confrontare più velocemente i segnali.
Il funzionamento di questo integrato può essere semplificato, se valutato in configurazione di open loop, schematizzandone il funzionamento con un transistor BJT a collettore aperto, in il collettore del transistor rappresenta l'output dell'opamp e l'emettitore è collegato a comune (Figura \ref{cir4:open_collector}).
Se colleghiamo anche una resistenza $R$ fra una tensione arbitraria $V_C$, la funzione di output dell'operazionale diventa

\begin{equation}
V_{out} = \bigg \{
\begin{array}{rl}
V_C & \mathrm{se} \quad V_+ < V_- \\
0 & \mathrm{se} \quad V_+ > V_- \\
\end{array}
\label{eq4:comparatore}
\end{equation}

dove $V_+$ e $V_-$ sono rispettivamente le tensioni all'ingresso non invertente ed invertente. Utilizzando la schematizzazione con il transistor, è come considerare la condizione di quest'ultimo nella seguente maniera

\begin{equation}
Q : \bigg \{
\begin{array}{rl}
\mathrm{interdizione} & \mathrm{se} \quad V_+ < V_- \\
\mathrm{saturazione} & \mathrm{se} \quad V_+ > V_- \\
\end{array}
\label{eq4:comparatore_Q}
\end{equation}

da cui è facile mostrare che la tensione di uscita è proprio (\ref{eq4:comparatore}).

\begin{wrapfigure}[16]{l}{0.5\textwidth}
  \begin{center}
    \includegraphics[width=0.280\textwidth]{../E04/latex/c_LM311.pdf}
  \end{center}
  \caption{Schema a collettore aperto dell'amplificatore operazionale LM311.}
  \label{cir4:open_collector}
\end{wrapfigure}

Infatti, quando Q è interdetto il generatore di tensione che fornisce $V_C$ non eroga corrente (l'impedenza dell'oscilloscopio è \SI{1}{\Mohm}) e non c'è caduta ai capi di $R$, dunque $V_{out}=V_C$; quando Q è in saturazione la tensione di uscita (tensione di collettore) si porta circa alla tensione di emettitore, dunque a zero\footnote{In realtà la tensione di uscita risulterà qualche centinaio di \si{\milli\volt} sopra lo $0$. Tale tensione serve a polarizzare il transistor affinché vada in saturazione oltre, ovviamente, alle cadute nelle altre parti della circuiteria interna dell'amplificatore.}. Si noti che la tensione di uscita può essere controllata solo grazie alla presenza della resistenza $R$.

Sfruttando questa funzione possiamo dunque confrontare la tensione in entrata su un ingresso con quella all'altro ingresso. Ad esempio, poniamo l'ingresso invertente ad una tensione di soglia $V_{ref}$: se vale la (\ref{eq4:comparatore}), ponendo un segnale in alternata sull'ingresso non invertente, l'operazionale restituirà una $V_{out}$ nulla se $V_{in}<V_{ref}$, altrimenti $V_C$. Valutiamo dunque il funzionamento del circuito ipotizzato, il cui schema circuitale è in Figura \ref{cir4:comparatore}, ponendo $V_{C}=5$\si{\volt}. Il grafico in Figura \ref{gr4:comparatore} mostra la tensione in entrata e in uscita.

\begin{figure}[ht]
 \centering
   {\includegraphics[width=5.5cm]{../E04/latex/c_comparatore.pdf}}
 \caption{Schema completo del comparatore utilizzato. La resistenza $R_p=(9.98 \pm 0.01)$ \si{\kilo\ohm}, come consigliato sul datasheet.}
 \label{cir4:comparatore}
\end{figure}

\begin{figure}[ht]
 \centering
   {\includegraphics[width=14.5cm]{../E04/latex/comp.pdf}}
 \caption{Grafico della tensione in entrata (onda sinusoidale con $V_{pp}=2$ \si{\volt} e $f=500$\si{\hertz}, tratteggiata) e in uscita (verde) in funzione del tempo. Notiamo che l'uscita rispetta (\ref{eq4:comparatore}).}
 \label{gr4:comparatore}
\end{figure}

%\begin{wrapfigure}[7]{l}{0.45\textwidth}
%  \begin{center}
%    \includegraphics[width=0.30\textwidth]{../E04/latex/c_comparatore.pdf}
%  \end{center}
%  \caption{Schema completo del comparatore utilizzato. La resistenza $R_p=(9.98 \pm 0.01)$ \si{\kilo\ohm}, come consigliato sul datasheet.}
%  \label{cir4:comparatore}
%\end{wrapfigure}
\subsection{Trigger di Schmidt non-invertente}

Quando nel segnale è presente un rumore è possibile che il valore di tensione restituito dal comparatore cambi più volte da $0$ a $V_{CC}$. Infatti, l'ampiezza in tensione del rumore può essere tale da passare sopra e sotto la tensione di soglia, invertendo la $V_{out}$ più volte secondo la \ref{eq4:comparatore}. Per ovviare a questo problema si sfrutta la retroazione positiva dell'amplificatore operazionale per riportare parte della $V_{out}$ sull'ingresso non invertente: ciò permette di creare un offset positivo sulla tensione all'ingresso non invertente quando $V_{in}>0$ e viceversa. Quindi, durante il confronto del segnale con la tensione di riferimento
\subsection{Oscillatore a rilassamento}
In questa parte dell'esperienza studieremo il funzionamento di un oscillatore a rilassamento. Tale circuito sfrutta l'instabilità della retroazione positiva per generare un segnale oscillante con una frequenza determinata dai componenti stessi del circuito. In Fig.(\ref{suca}) è riportato lo schema circuitale. 

L'elemento circuitale fondamentale per il funzionamento dell'oscillatore a rilassamento è il condensatore C. Per comodità a $t=0$ assumiamolo scarico e $V_{out}=+V_{sat}$. Ovviamente, per effetto della retroazione negativa, la tensione all'ingresso invertente tenderà ad aumentare ($V_{inv}$ è determinata da quanta carica è accumulata sul condensatore). Non appena $|V_{inv}|>|V_{ninv}|$, avremo uno switch della tensione in uscita, ovvero $+V_{sat} \rightarrow -V_{sat}$. Il condensatore inizierà dunque a scaricarsi. Successivamente avremo un altro swich della tensione in uscita e il ciclo si ripeterà. Il periodo di tale oscillazione può essere calcolato analiticamente risolvendo il circuito. Il valore che si ottiene è $T=2RCln(1+\frac{2R_1}{R_2})$.  

Nei seguenti grafici sono riportati i risultati ottenuti per valori di $R=(9.99\pm0.01)\si{\kilo\ohm}$ e $R=(99.93\pm0.01)\si{\kilo\ohm}$.


$$GRAFICO$$

$$GRAFICO$$

Nella seguente tabella riportiamo i valori teorici e sperimentali di periodo dell'oscillatore a rilassamento.

\begin{tabular}{|l|l|l|}
\hline
$R [\si{\kilo\ohm}]$	&  $T_{exp} [\si{\milli\second}]$          & $T_{teo} [\si{\milli\second}]$       \\
\hline
$9.99\pm0.01 $ & $2.22\pm0.01$ & \\
\hline
$99.93\pm0.01 $ & $22.16\pm0.01$ & \\
\hline
\end{tabular} 

\subsection{Interruttore crepuscolare}

Nell'ultima parte dell'esperienza abbiamo assemblato un interruttore crepuscolare, cioè un circuito elettronico che alimenta un carico (LED) in condizioni di bassa luminosità ambientale e interrompe l'alimentazione in condizioni di alta luminosità ambientale.

\begin{figure}[ht]
 \centering
   {\includegraphics[width=12.5cm]{../E04/latex/c_crepuscolare.pdf}}
 \caption{DIDASCALIA!! $R_F = \SI{1}{\Mohm}$}
 \label{gr4:crepuscolare}
\end{figure}

Un circuito di questo tipo necessità di almeno due componenti: il sensore di luminosità (fototransistor OP550) che fornisce la variabile circuitale collegata alla luminosità ambientale e un blocco di comparazione che \textit{decide} se alimentare o meno il carico confrontando il segnale dato dal primo blocco con un segnale di riferimento.

Il componente principale del primo blocco è il fototransistor OP550 che agisce come una sorgente di corrente, generando una corrente dell'ordine del \si{\uA} se esposto alla luce.
Tale corrente è però difficilmente manipolabile, pertanto abbiamo deciso di utilizzare un convertitore corrente-tensione per trasformare e amplificare il segnale in modo tale da poter essere utilizzata come segnale di ingresso del secondo blocco.

Un convertitore corrente-tensione può essere implementato facilmente utilizzando un opamp $\mu$A741 in modalità di amplificatore invertente con retroazione negativa (vedi Figura \ref{cir4:crepuscolare}).
In questo modo, la tensione in uscita dell'opamp sarà legata alla corrente che attraversa il sensore dalla seguente equazione:

\begin{equation}
	V_{out} = R_F \,\, I_{OP550}
\end{equation}
dove $R_F$ è la resistenza posta sul ramo di retroazione (nel nostro caso $R_F = \SI{1}{\Mohm}$) e $I_{OP550}$ è il modulo della corrente passante nel fototransistor.

Il secondo blocco del circuito è invece composto da un opamp LM311 utilizzato come comparatore con isteresi (o trigger di Schmidt).
Sul ramo di retroazione e tra il terminale positivo e l'ingresso del segnale - che corrisponde all'uscita dell'opamp $\mu$A741 - sono state poste due resistenze per creare l'isteresi affinché eventuali rumori di fondo non disturbassero il funzionamento del circuito.
Nel nostro caso abbiamo utilizzato due resistenze $R_1 =$ \SI{5}{\kohm} e $R_2=$ \SI{100}{\kohm} in modo da avere un ampio margine di sicurezza.

Come si può notare dal grafico in Figura \ref{gr4:neon} il rumore ambientale dato dalle lampade al neon ha una forma d'onda sinusoidale di frequenza \SI{100}{\hertz} e ampiezza picco-picco \SI{80}{\mV}, pertanto sarebbe stata sufficiente un'isteresi di appena $\frac{V_{CC}}{100} \approx \SI{150}{\mV}$

Per impostare la tensione di soglia, invece, abbiamo creato un partitore collegando i due connettori esterni di un trimmer resistivo da \SI{10}{\kohm} a +\SI{15}{\V} e a comune, mentre abbiamo collegato il terminale centrale del trimmer all'ingresso invertente del LM311.
In questo modo potevamo scegliere la tensione di soglia girando semplicemente un cacciavite.

Infine abbiamo posto il carico da alimentare tra l'alimentazione positiva +\SI{15}{\V} e la resistenza di pull-up ($R_P=$ \SI{1}{\kohm}).

\begin{figure}[ht]
 \centering
   {\includegraphics[width=14.5cm]{../E04/latex/neon.pdf}}
 \caption{Rumore ambientale determinato dalla luce delle lampade al neon del laboratorio incidente sul fototransistor e amplificato dall'opamp $\mu$A741. I dati sono stati ricavati separando i due blocchi del circuito e misurando la tensione all'uscita dell'opamp $\mu$A741.}
 \label{gr4:neon}
\end{figure}
\subsection*{Conclusioni}

In questa esperienza abbiamo analizzato il funzionamento di circuiti comparatori, molto utili quando vogliamo confrontare due segnali ottenendo un semplice valore digitale (on-off). Tuttavia, vista la sensibilità degli op-amp, abbiamo bisogno di un modo per evitare che il rumore influenzi le nostre comparazioni. Per far questo abbiamo introdotto i circuiti con isteresi, ovvero con feedback-positivo. \'E stato poi realizzato un interruttore crepuscolare. Questi interruttori, contenenti un fototransistor, permettono ad esempio l'accensione (e spegnimento) delle luci stradali quando tramonta il sole. Non si vuole però che una semplice nuvola passeggera faccia accendere le luci. Bisogna dunque introdurre un' insteresi per far sì che il rumore (la nuvola) non faccia scattare l'interruttore. Infine, abbiamo costruito un oscillatore a rilassamento. Peculiarità di tale circuito è che esso oscilla tra $+V_{sat}$ e $-V_{sat}$ sebbene sia alimentato con tensioni continue. La frequenza di tale oscillatore è univocamente determinata dai valori delle componenti circuitali.


\clearpage
\section{14.10.2014 - Raddrizzatore, Amplificatori e Termoresistenza}

In questa esperienza progetteremo un raddrizzatore di precisione (utilizzabile anche per segnali di bassa tensione). Dopodiché valuteremo la capacità di un amplificatore differenziale e dell'integrato $AD622$ di abbattere i guadagni di modo comune. Infine utilizzeremo una termoresistenza per ricavare la temperatura in alcune situazioni.

\subsection*{Strumenti e materiali}

\begin{itemize} [noitemsep]
\item Oscilloscopio Agilent DSO-X 2002A (bandwidth \SI{70}{\mega\hertz}, sample rate \num{2} GSa/s);
\item Generatore di tensione continua Agilent E3631A (max $\pm \, \SI{25}{\volt}$ o $\pm \, \SI{6}{\volt}$);
\item Generatore di forme d'onta Agilent 33120A con range di frequenza da \SI{100}{\micro\hertz} a \SI{15}{\mega\hertz};
\item Multimetro Agilent 34410A a sei cifre e mezza;
\item Un amplificatore operazionale $\mu$A741;
\item Un integrato AD622;
\item Una termoresistenza PT100;
\item Resistenze e capacità di vari valori;
\item Due trimmer multigiro da \SI{10}{\kilo\ohm};
\item Breadboard e cablaggi vari.
\end{itemize}

\subsection{Raddrizzatore di precisione}

\subsubsection{Premessa sui raddrizzatori}

I raddrizzatori sono dei circuiti il cui scopo è quello di effettuare il modulo di un segnale in ingresso preservandone la forma. In precedenti corsi di optoelettronica abbiamo incontrato alcuni semplici raddrizzatori i cui elementi circuitali chiave erano i diodi. Questi dispositivi sono in grado di portarsi in conduzione, e quindi di chiudere il ramo di circuito in cui si è posto il diodo, se la tensione al polo positivo è maggiore di quella al polo negativo; altrimenti il diodo apre il circuito. In realtà, ciò accade unicamente nel caso in cui il diodo sia ideale: per polarizzare un diodo è infatti necessario creare una differenza di tensione ai suoi capi, detta tensione di soglia $V_d$, che ha un valore di $\approx 0.6$ \si{\volt}; altrimenti il diodo risulterà come non polarizzato e dunque non in conduzione. Questo fatto, nella creazione di raddrizzatori può creare alcuni svantaggi, soprattutto se necessaria un'alta precisione (lavorando con piccoli valori di tensione): il segnale in uscita sarà con la stessa forma d'onda del segnale in entrata, ma con un offset di $-0.6$\si{\volt} dato dalla $V_d$; inoltre, se i segnali in entrata hanno valori più piccoli della tensione di soglia, il diodo non entrerà mai in conduzione. Di seguito presentiamo alcuni circuiti che risolvono in parte queste problematiche.

\subsubsection{Raddrizzatore con un diodo}

Analizziamo il circuito in Figura ?????. Vogliamo dimostrare che, qualunque sia il segnale in entrata, il contributo alla tensione di uscita dato dalla tensione di soglia del diodo è trascurabile.

Innanzitutto notiamo che, dato l'alto guadagno dell'operazionale, la tensione in entrata per polarizzare il diodo è un valore molto basso (con un guadagno di $2\times 10^5$ la tensione necessaria è data da $0.6 \si{\volt}/2\times 10^5 \approx 3 \si{\micro\volt}$); dunque non dobbiamo preoccuparci che il segnale in entrata debba avere un valore minimo per far funzionare il circuito.

Vale per l'operazionale l'equazione (\ref{eq3:regola_opamp})\footnote{In realtà bisognerebbe considerare anche la dipendenza dalla frequenza. Abbiamo però notato che, ben prima che si possa apprezzare l'abbattimento dell'amplificazione, incorreremo in un altro problema dato dalla velocità di polarizzazione dei transistor all'interno dell'operazionale che ci porterà a scartare questo circuito come raddrizzatore di precisione ad alte frequenze.}
\begin{equation}
V_{ol}=A (V^+-V^-)
\label{eq5:regola_opamp_NOFREQDEP}
\end{equation}
dove $V_{ol}$ è la tensione all'uscita dell'amplificatore operazionale.
Valutiamo ora il circuito in due casi: $V_{in}<0 e V_{in}>0$\footnote{Grazie alle considerazioni sopra possiamo trascurare la tensione in ingresso necessaria per polarizzare il diodo}.

\paragraph*{Caso $V_{in}<0$}

Nel primo caso, l'OPAMP amplifica il segnale negativo in entrata sull'uscita rispettando la (\ref{eq5:regola_opamp_NOFREQDEP}) con $V_-=0$: infatti, essendo il diodo interdetto (il polo positivo è più negativo di quello negativo), non passa alcuna corrente per $R_L$ rendendo $V^-=V_{out}=0$.

\paragraph*{Caso $V_{in}>0$}

Nel secondo caso invece il diodo è in conduzione. Consideriamo dunque, come si nota dalla Figura ??????
$$V^+=V_{in} \qquad V^-=V_{out}$$
Dato che un diodo in conduzione ha una caduta in buona approssimazione costante $V_d$, vale anche che
$$V_{ol}=V_{out}+V_d$$
Sostituendo questi risultati in (\ref{eq5:regola_opamp_NOFREQDEP}) otteniamo le equazioni cercate
\begin{equation}
V_{out}=\frac{1}{1+A} (A V_{in} - V_d) \approx V_{in} \qquad V_{ol} \approx V_{in}+V_d
\label{eq5:leggi_1.1}
\end{equation}

Di seguito proponiamo i grafici delle leggi sopra ricavate plottate con i dati sperimentali.

Dal plot del segnale in entrata ed in uscita in funzione del tempo alla frequenza di ?????? \si{\hertz} (grafico in Figura ???????), possiamo notare un troncamento della forma d'onda del segnale in uscita. Abbiamo ipotizzato che tale fenomeno fosse dovuto alla limitata velocità che i transistor all'interno dell'operazionale hanno di passare dalla saturazione all'interdizione; tale fenomeno è anche amplificato dal fatto che, quando il diodo è interdetto, $V_{ol} \approx -V_{CC}$: il segnale deve quindi passare da una tensione di saturazione negativa alla tensione necessaria per raddrizzare il segnale in uscita in un tempo troppo breve, e ad alte frequenze ciò risulta in un troncamento. Per risolvere tale problema utilizziamo il circuito esposto al sottoparagrafo successivo.

\subsubsection{Raddrizzatore con due diodi}

Per risolvere il problema visto in precedenza, utilizziamo un secondo diodo posto come nel in Figura ?????? e un circuito di retroazione controllata in guadagno. Anche in questo caso dobbiamo valutare la risposta del circuito in casi diversi: $V_{in}=0$, $V_{in}>0$ e $V_{in}<0$.

\paragraph*{Caso $V_{in}=0$}

Nel primo caso abbiamo entrambi i diodi interdetti. Infatti, se la tensione in ingresso è nulla non vi è la tensione necessaria a polarizzare i diodi, che rimangono non polarizzati. Dunque $V_{out}=V_{ol}=0$.

\paragraph*{Caso $V_{in}>0$}

Tenendo conto della (\ref{eq5:regola_opamp_NOFREQDEP}), considerando che $V_{in}$ viene portata sull'ingresso invertente, la tensione su $V_{ol}$ sarà negativa e quindi più bassa di $V_{out}$: dunque $D_2$ sarà interdetto. Allo stesso tempo il polo positivo di $D_1$ è più positivo (a comune virtuale) di quello negativo ($V_{ol}$), dunque tale diodo risulta in conduzione. Inoltre, $V_{out}$ risulta nulla in quanto sul ramo di circuito identificato da $R_2$ e $R_L$ non passa corrente (sono collegate rispettivamente a comune virtuale e comune). $V_{ol}$ sarà invece fissata dalla tensione di soglia del diodo, cioè $V_{ol}=-V_d$.

\paragraph*{Caso $V_{in}<0$}

Con considerazioni analoghe al caso precedente, possiamo concludere che $D_1$ risulta interdetto, mentre $D_2$ in conduzione. Quindi il circuito è un amplificatore invertente di segnale con guadagno $-R_2/R_1$.

Valutiamo il contributo del diodo nella tensione di uscita del circuito $V_{out}$ e dell'operazionale $V_{ol}$. Consideriamo le correnti nel nodo $V_N$
$$\frac{V_{in}-V_N}{R_1} + \frac{V_{out}-V_N}{R_2} = 0$$
Osservando che $V_N$ è a ground virtuale, che $V_{out}=V_{ol}-V_d$, e considerando l'operazionale come ideale e che vale che $V_{out}=V_{ol}-V_d$, dalla (\ref{eq5:regola_opamp_NOFREQDEP}) otteniamo
$$V_{out}=-\frac{R_2}{R_1}V_{in} \qquad \frac{R_2}{R_1} V_{in} + V_d$$

Di seguito presentiamo i grafici delle tensioni.

INSERIRE GRAFICI E COMMENTARE!
\subsection{Amplificatore differenziale}

\begin{wrapfigure}[14]{l}{0.5\textwidth}
  \begin{center}
    \includegraphics[width=0.280\textwidth]{../E05/latex/c_teo_diff_amp.pdf}
  \end{center}
  \caption{...}
  \label{cir5:diff_amp_teo}
\end{wrapfigure}

In questa parte dell'esperienza cercheremo di capire e analizzare un amplificatore differenziale (\ref{cir5:diff_amp}). Oltre al fatto che scegliendo i valori di resistenza possiamo decidere il guadagno, la proprietà più importante è sicuramente il vantaggio che si ha rispetto al rumore in modo comune. Per capirne il motivo analizziamo il circuito riportato in figura (\ref{cir5:diff_amp_teo}).

Sfruttiamo il principio di sovrapposizione per valutare la tensione in uscita in funzione delle due di ingresso. Ricordiamo che la sovrapposizione può essere applicata per sistemi lineari. Ciò è molto utile in quanto ci permette di separare l'analisi circuitale e ricondurci a casi più semplici.

Chiamiamo $V_-$ la tensione all'ingresso invertente e $V_+$ quella all'ingresso non invertente. Poniamo $V_1=0$. La tensione le punto $V_+$ sarà ovviamente data dal partitore con $R$ ed $mR$. Assumendo l'op-amp ideale, abbiamo $V_+=V_-$. Possiamo dunque scrivere:

\begin{equation}
\begin{cases} V_+=V_-  \\ V_+=V_2\frac{mR}{R+mR} \\ V_-= V_{out2}\frac{R}{R+mR}\end{cases}  \Rightarrow V_{out2}=mV_2
\label{eq 5: vout2}
\end{equation}

\begin{wrapfigure}[13]{r}{0.5\textwidth}
  \begin{center}
    \includegraphics[width=0.280\textwidth]{../E05/latex/c_diff_amp.pdf}
  \end{center}
  \caption{...}
  \label{cir5:diff_amp}
\end{wrapfigure}

Analogamente, poniamo $V_2=0$. Otteniamo le seguenti equazioni del circuito:

\begin{equation}
\begin{cases} V_+=V_-=0  \\ \frac{V_1}{R}+\frac{V_{out2}}{mR}=0 \end{cases}  \Rightarrow V_{out1}=-mV_1
\label{eq 5: vout1}
\end{equation}

Sommando ora (\ref{eq 5: vout2}) e (\ref{eq 5: vout2}) otteniamo trivialmente $V_{out}=m(V_2-V_1)$.

Possiamo ora passare ad analizzare il circuito da cui siamo partiti, ovvero quello affetto da noise riportato in figura (\ref{cir5:diff_amp}). Possiamo scrivere, senza perdere di generalità,
$$
\begin{cases} V_1=V_{in}+V_{noise} \\ V_2=V_{noise} \end{cases}  \Rightarrow V_{out}=m(V_{noise}-(V_{in}+V_{noise}))=-V_{in}
$$

L'amplificatore differenziale ci permette dunque di eliminare in modo efficace il rumore di modo comune. 

Come vediamo, la resistenza collegata tra comune e $V_+$ è in realtà un trimmer. Ciò è necessario in quanto, non essendo esattamente uguali le resistenze ed essendo reale l'op-amp, prima di utilizzare il circuito è necessario bilanciare il circuito in modo da ottenere un segnale il più piccolo possibile quando i segnali in ingresso sono uguali. 

Il laboratorio abbiamo utilizzato come sorgente di noise il generatore di forme d'onda e come $V_{in}$ il generatore di tensione costante. Abbiamo utilizzato delle R da \SI{10}{\kilo\ohm} e $mR=2R$. La resistenza variabile è stata costruita mettendo in serie una da \SI{10}{\kilo\ohm} con un trimmer sempre da \SI{10}{\kilo\ohm}. Per controllare il bilanciamento abbiamo connesso entrambi gli ingressi al generatore di forme d'onda ed è stato utilizzato un segnale sinusoidale di $20Vpp$. Il segnale in uscita è stato dunque riportato sull'oscilloscopio. Abbiamo notato che cambiando il valore della resistenza con il trimmer il segnale in uscita cambiava. Il miglior bilanciamento che siamo riusciti a raggiungere a portato la tensione picco-picco del rumore a \SI{0.5}{\milli\volt}. Abbiamo dunque riattaccato il generatore di tensione costante e alimentato il circuito con $-2Vpp$. Il rumore è sempre stato simulato con una sinusoidale di $20Vpp$. I risultati sono riportati nella seguente figura.

\begin{figure}[ht]
 \centering
   {\includegraphics[width=0.75\textwidth]{../E05/latex/amp_diff.pdf}}
 \caption{Come vediamo dal grafico, il rumore che abbiamo simulato ha un'ampiezza di 20 Volt picco-picco (sinusoide in nero). Nonostante ciò, l'uscita (in verde) non è affatto influenzata e risulta amplificata di 2 volte, come stimato precedentemente dai calcoli teorici. }
 \label{gr5:amp_diff}
\end{figure}

Come vediamo, il rumore è stato completamente eliminato dall'amplificatore differenziale. Abbiamo successivamente provato a sbilanciare il circuito, per vedere l'effetto del rumore sul segnale in uscita. E' stato dunque variata la resistenza con trimmer e abbiamo utilizzato un rumore con dei picchi molto accentuati, così da vedere bene gli effetti sul segnale in uscita. 

\begin{figure}[ht]
 \centering
   {\includegraphics[width=0.75\textwidth]{../E05/latex/sbil_amp_diff.pdf}}
 \caption{Dopo lo sbilanciamento del circuito non è più garantita la soppressione del rumore in modo comune. Come vediamo, il segnale in uscita risulta distorto. E' dunque necessario calibrare bene il circuito prima di effettuare esperimenti nei quali si vuole eliminare un rumore comune ad entrambi i segnali in ingresso.}
 \label{gr5:sbil_amp_diff}
\end{figure}

Come vediamo, il segnale risulta disturbato dal rumore. Inoltre, il guadagno può essere cambiato solo se si cambiano 2 resistenze nel circuito. Una soluzione efficace che inoltre rimedia alla bassa impedenza in ingresso è quella di utilizzare un Amplificatore per Strumentazione.

\subsection{Instrumentation Amplifier}

Il circuito visto al paragrafo precedente presenta dei problemi. Per prima cosa l'impedenza in ingresso vista da un eventuale generatore di segnale è bassa, e ciò rende necessaria una potenza maggiore da parte del generatore per mantenere la tensione ai suoi capi; il secondo luogo, il guadagno non è impostabile (se non cambiando le resistenze ad ogni applicazione: procedura scomoda nelle applicazioni pratiche).

Per ovviare a questi problemi, valutiamo il circuito in Figura ??????. Si nota subito che l'impedenza in ingresso è molto alta, in quanto gli amplificatori operazionali hanno un valore di impedenza molto alta ai loro ingressi (non a caso sono usati in configurazione follower per adattare le impedenze).

Verifichiamo invece che la resistenza posta in mezzo varia effettivamente il guadagno del circuito. Dividiamo il circuito in due parti, calcolandoci prima la tensione fra il punto A e B\footnote{Si noti che un circuito che abbia come $V_{out}$ la differenza fra questi due punti sarebbe flottante sull'uscita. Infatti, una eventuale resistenza di carico messa fra il punto A e B non avrebbe un riferimento di comune.}, per poi la tensione di uscita dell'intero circuito. Per la prima parte, considerando i due operazionali ideali, otteniamo che le tensione agli ingressi invertenti sono $V_{OP_1}^- = V_1$ e $V_{OP_2}^- = V_2$. Dunque ai capi di $R_g$ è presente una differenza di potenziale $V_2-V_1 = \Delta V_{in}$; per l'idealità degli OPAMP, abbiamo inoltre che la corrente che passa per le resistenze $R_1$ ed $R_g$ sono uguali. Otteniamo dunque, sommando le cadute di potenziale
$$\Delta V_{AB} = \frac{\Delta V_{in}}{R_g} R_g + \frac{\Delta V_{in}}{R_g} R_1 + \frac{\Delta V_{in}}{R_g} R_2$$
cioè
\begin{equation}
\Delta V_{AB} = \Delta V_{in} \left(1+\frac{2R_1}{R_g}\right)
\label{eq5:TEMP_calcoli}
\end{equation}

Valutiamo ora la seconda parte del circuito data da $OP_3$. Per quanto riguarda la tensione all'ingresso non invertente, questa è data da un semplice partitore (non scorre corrente nell'ingresso dell'opamp)
$$V^+=V_B \frac{R_3}{R_3+R_2}$$
All'ingresso non invertente vale invece la legge di Kirkhhoff per i nodi
$$\frac{V_A-V^-}{R_2}+\frac{V_{out}-V^-}{R_3}=0$$
da cui
$$V^-=V_A \frac{R_3}{R_2+R_3} + V_{out} \frac{R_2}{R_2+R_3}$$

Uguagliando $V^-$ e $V^+$ (OPAMP ideale), tenendo conto della (\ref{eq5:TEMP_calcoli}), otteniamo dunque
$$V_{out}=-\Delta V_{in} \left(1+\frac{2R_1}{R_g}\right)\frac{R_3}{R_2}$$
Modificando la resistenza $R_g$ possiamo dunque controllare il guadagno del circuito.

In realtà, nell'esperienza abbiamo utilizzato un circuito integrato (più precisamente l'AD622) che ha come qualità il fatto di avere un'alta precisione sul valore delle resistenze (abbiamo visto nel precedente circuito che uno squilibrio anche minimo fra i valori di resistenze che dovrebbero essere uguali può portare ad un guadagno diverso da quello teorico). Il suo guadagno è dato dall'equazione, fornita dal costruttore
$$G=1+\frac{50.5 \si{\kilo\ohm}}{R_g}$$


\subsection{Misure di temperatura}

\begin{wrapfigure}[14]{l}{0.45\textwidth}
\centering
\includegraphics[width=.25\textwidth]{../E05/latex/c_PT100_2wire.pdf}
\caption{Circuito di controllo della temperatura a distanza: PT100 è il termoresistore al platino, l'ohmetro misura la resistenza in funzione della temperatura e le due resistenze da \SI{10}{\ohm} simulano i cavi lunghi.}
\label{cir5:2wire}
\end{wrapfigure}

Nell'ultima parte dell'esperienza abbiamo eseguito misure di temperatura utilizzando una termoresistenza al platino PT100.
Essa non è altro che una resistenza di cui sono noti con sufficiente precisione alcuni parametri: il valore ad una data temperatura e il coefficiente di temperatura $\alpha$.

Sapendo che il valore di resistenza è dipendente dalla resistività elettrica del materiale $\rho$ che a sua volta dipende dalla temperatura (nel caso di metalli), secondo le seguenti relazioni:
	$$	R = \frac{\rho L}{S}
			\qquad \qquad
		\rho = \rho_0 \left[ 1 + \alpha \left( T - T_0 \right) \right]$$
dove L è la lunghezza del materiale e S la sezione, mentre $\rho$ dipende dalla resistività $\rho_0$ alla temperatura di riferimento $T_0$ e dal coefficiente di temperatura $\alpha$.
Combinando le due equazioni precedenti si ottiene la dipendenza del valore di resistenza dalla temperatura:
\begin{equation}
R = R_0 \left[ 1 + \alpha \left( T - T_0 \right) \right]
\end{equation}

\begin{wrapfigure}[17]{r}{0.45\textwidth}
\centering
\includegraphics[width=.3\textwidth]{../E05/latex/c_PT100_4wire.pdf}
\caption{Configurazione di misura a quattro fili: nella maglia esterna scorre la corrente data dal generatore mentre nella maglia interna non scorre corrente. Il voltmetro ha un'impedenza idealmente infinita.}
\label{cir5:4wire}
\end{wrapfigure}

Invertendo la formula si ottiene:
\begin{equation}
T = T_0 + \frac{1}{\alpha}\left( \frac{R}{R_0}-1 \right)
\end{equation}
Nel nostro caso il costruttore ci ha fornito i dati relativi alla temperatura di riferimento $T_0 =$ \SI{0}{\degreeCelsius}:
$$R_0 = \SI{100}{\ohm} \quad \quad \alpha = \SI[per-mode=symbol]{0.3850}{\ohm\per\degreeCelsius}$$

Conoscendo il comportamento della termoresistenza abbiamo simulato il caso in cui si volesse controllare la temperatura di un laboratorio a distanza.
I circuito che simulava questa situazione è riportato in Figura \ref{cir5:2wire}

Abbiamo notato subito che le misure ricavate con questo circuito soffrivano dell'errore sistematico dato dalle resistenze in serie da \SI{10}{\ohm}.
Infatti il valore di resistenza letto dall'ohmetro era di \SI{130}{\ohm}, che dava in risultato una temperatura ambiente di $\approx$ \SI{78}{\degreeCelsius}, evidentemente errata.

Per evitare questo tipo di problemi abbiamo effettuato una misura a quattro fili.
In questa modalità il multimetro funzionava come generatore di corrente costante tramite due boccole e da voltmetro con le altre due boccole utilizzate.
Lo schema riportato in Figura \ref{cir5:4wire} mostra come la misura questa volta sia effettuata acquisendo la differenza di potenziale ai capi della termoresistenza data dalla corrente del generatore.
In questa configurazione non vi è l'errore sistematico della misura a due fili, in quanto grazie all'alta impedenza del voltmetro, nella maglia interna non scorre corrente e pertanto ai capi delle resistenze su tale maglia non vi è alcuna caduta di potenziale.

Nella tabella seguente sono riportati i valori per $T_0 =$ 0, per due temperature vicine alla temperatura ambiente, per una temperatura ottenuta raffreddando con una bomboletta di refrigerante e riscaldando il termoresistore a temperatura corporea.

\begin{center}
{\renewcommand{\arraystretch}{1.4}%
\begin{tabular}{c|c c c c c}
$R$ [\si{\ohm}] 		& $R_0 =$ 100 	& 110.24 & 111.32 & 83.37 & 114.71\\ 
\hline 
$T$ [\si{\degreeCelsius}] 	& $T_0 =$ 0 	& 26.6 & 29.4 & -43.2 & 38.2 \\ 
\end{tabular}}
\end{center}
\subsection*{Conclusioni}
In questa esperienza abbiamo studiato alcuni circuiti utili a trasformare un segnale senza modificare l'informazione in esso contenuta: in particolare abbiamo esaminato circuiti raddrizzatori che non diminuiscono l'ampiezza del segnale come accade con i circuiti passivi quali i semplici diodi o il ponte raddrizzatore (ponte di Graetz).

Successivamente abbiamo studiato il cosiddetto amplificatore da strumentazione, in un circuito didattico costruito con un $\mu$A741 le cui limitazioni sono date dal guadagno fisso e dalla bassa impedenza in entrata e integrato in un dual inline package (AD622).
Ne abbiamo testato il funzionamento misurando lo sbilanciamento di un ponte di Wheatstone.

Infine abbiamo studiato una problematica di una misura a distanza data dalla lunghezza dei cavi.
Abbiamo risolto il problema con una misura a quattro fili.

\clearpage
\section{22.10.2014 - Termostatazione}

In questa esperienza costruiremo un termometro elettronico utilizzando una Pt100. Realizzeremo poi un sistema di controllo proporzionale di temperatura utilizzando il termometro da noi costruito. 

\subsection*{Strumenti e materiali}

\begin{figure}[htc]
\begin{itemize} [noitemsep]
	\item Generatore di forme d'onta Agilent 33120A con range di frequenza da \SI{100}{\micro\hertz} a \SI{15}{\mega\hertz};
	\item Oscilloscopio Agilent DSO-X 2002A (bandwidth \SI{70}{\mega\hertz}, sample rate \num{2} GSa/s);\newline
	\begin{minipage}{0.65\textwidth}
		\vspace{0.4mm}
%		\begin{itemize} [noitemsep]
%		\item Oscilloscopio Agilent DSO-X 2002A (bandwidth \SI{70}{\mega\hertz}, sample rate \num{2} GSa/s);
		\item Generatore di tensione continua Agilent E3631A (max $\pm \, \SI{25}{\volt}$ o $\pm \, \SI{6}{\volt}$);
%		\item Generatore di forme d'onta Agilent 33120A con range di frequenza da \SI{100}{\micro\hertz} a \SI{15}{\mega\hertz};
		\item Multimetro Agilent 34410A a sei cifre e mezza;
		\item Tre amplificatori operazionali OP07;
		\item Un integrato AD622 e un REF02;
		%\item Un integrato REF02;
		\item Una termoresistenza PT100;
		\item Transistor di potenza NPN 2N2222;
		\item Resistenze e capacità di vari valori;
		\item Trimmer multigiro da \SI{10}{\kilo\ohm} e \SI{1}{\kilo\ohm};
		\item Breadboard e cablaggi vari.
%		\end{itemize}
	\end{minipage}
	\begin{minipage}{0.3\textwidth}
			\centering
			\includegraphics[width=3.5cm]{../E06/latex/REF02.pdf}
			\caption{Piedinatura dell'integrato\newline REF02.}
			\label{cir6:REF02}
%2 è alimentato con +\SI{15}{\volt}, 4 è collegato a ground, 6 restituisce una tensione costante di +\SI{5}{\volt}.}
	\end{minipage}
\end{itemize}
\end{figure}
\vspace{-.8cm}

\subsection{PRIMO PARAGRAFO DEL GRUPPO AVANZATO}
\subsection{SECONDO PARAGRAFO MEGLIO DEL PRIMO E DELL'INTRODUZIONE DEL GRUPPO AVANZATO}
%\subsection{TERZO PARAGRAFO MEGLIO DEL PRIMO E DELL'INTRODUZIONE E ANCHE UN PO' PIU' COLORATO DEL GRUPPO AVANZATO}
%\subsection{QUARTO PARAGRAFO MEGLIO DEL PRIMO E DELL'INTRODUZIONE E ANCHE UN PO' PIU' COLORATO E SGAMATO DEL GRUPPO AVANZATO}
\subsection*{Conclusioni}
COMMENTARE il commento per caso a casa di casa con casa a caso in conclusione.


\clearpage
\section{28.10.2014 - Circuito per elettrocardiogramma (ECG)}

In questa esperienza ingegneristica costruiremo un circuito per effettuare un elettrocardiogramma (ECG).

\subsection*{Strumenti e materiali}

\begin{itemize} [noitemsep]
	\item Oscilloscopio Agilent DSO-X 2002A (bandwidth \SI{70}{\mega\hertz}, sample rate \num{2} GSa/s);
		\item Generatore di tensione continua Agilent E3631A (max $\pm \, \SI{25}{\volt}$ o $\pm \, \SI{6}{\volt}$);
		\item Multimetro Agilent 34410A a sei cifre e mezza;
		\item Due pile da \SI{9}{\volt};
		\item Due elettrodi;
		\item Amplificatori operazionali OP07;
		\item Un integrato AD622 e un ISO124;
		\item Resistenze e capacità di vari valori;
		\item Breadboard e cablaggi vari;
		\item Paziente (umano).
\end{itemize}

\begin{wrapfigure}[20]{r}{0.24\textwidth}
\centering
\includegraphics[width=.18\textwidth]{../E07/latex/mustache.pdf}%human.pdf}%doge.pdf}
\caption{Connessioni al corpo poco umano di uno degli sperimentatori.}
\label{fig7:human}
\end{wrapfigure}

\subsection*{Dati di progetto}

\begin{wrapfigure}[20]{r}{0.24\textwidth}
\centering
\includegraphics[width=.18\textwidth]{../E07/latex/human.pdf}
\caption{Connessioni al corpo umano}
\label{fig7:human}
\end{wrapfigure}

Per progettare un circuito che legga il segnale elettrico del nostro battito cardiaco a livello della cute, dobbiamo avere innanzitutto ben presenti alcuni punti fondamentali:
\begin{itemize} [noitemsep]
	\item Sulla superficie del corpo il segnale ECG è normalmente di circa \SI{1}{mV} (al massimo raggiunge i \SI{4}{mV}) ed ha una frequenza di circa \SI{1.25}{\Hz}.
	\item Gli elettrodi a nostra disposizione, che realizzano il collegamento tra la strumentazione e il corpo, creano un potenziale di circa \SI{700}{\mV}, quindi di circa \num{3} ordini di grandezza maggiore del segnale cardiaco.
	\item Il segnale da misurare è piccolo, quindi dobbiamo porre particolare attenzione alle fonti di rumore elettromagnetico dati da effetti capacitivi e/o induttivi. Inoltre dobbiamo evitare che i cavi e il paziente si muovano durante la misurazione.
	\item È necessario che il paziente sia isolato dal lato di elaborazione del segnale, alimentato dall'alimentatore da banco.
\end{itemize}

Analizziamo ora le parti che formeranno il circuito completo.

\subsection{Amplificatore differenziale}

Questo primo stadio in entrata (operazionale U1 in Figura \ref{cir8:compensation}) è necessario per avere un'alta impedenza in ingresso (e dunque un valore maggiormente fedele alla reale ddp del corpo), ma soprattutto per eliminare i segnali di modo comune generati dagli elettrodi. In testa ed in coda a questo stadio sono però presenti delle capacità che creano due filtri: rispettivamente un passa basso ed un passa alto.

Il primo filtro passa basso posto in testa ad U1, e composto dalle capacità $C_c$ e $C_d$, è necessario per abbattere eventuali rumori ambientali causati dalle radio-frequenze. Date le capacità in Figura \ref{cir8:compensation}, il datasheet dell'operazionale ci fornisce le seguenti frequenze di taglio di modo differenziale
\begin{equation*}
	\frac{1}{2 \pi R ( C_d + C_c ) } = \SI{194}{\Hz}
\end{equation*}
e di modo comune
\begin{equation*}
	\frac{1}{2 \pi R C_c} = \SI{4}{\kHz}
\end{equation*}

Il secondo filtro passa alto ha invece con frequenza di taglio pari a circa \SI{3}{\Hz}. Essendo posto all'uscita dell'operazionale (capacità $C_1$ in Figura \ref{cir8:compensation}, elimina eventuali segnali continui in output.

\subsection{Active gards}
Polarizzo la calza del cavo coassiale che connette l’elettrodo all’amplificatore ad un
potenziale pari alla tensione di modo comune.
In tal modo riduco l’effetto delle capacità parassite, e aumento la reiezione di modo
comune.
\begin{itemize}
\item il cavo che porta il segnale “vede” una differenza di potenziale molto bassa. Limito le
perdite nell’isolante
\item entrambi i cavi “vedono” lo stesso potenziale, indipendentemente dalla loro posizione sul
banco)
\end{itemize}

\subsection{Amplificatore di isolamento}
Isola galvanicamente i 2 lati del circuito: quello collegato al paziente ed alimentato a
batteria e quello di uscita ed elaborazione del segnale alimentato dall’alimentatore da
banco.
\begin{itemize}
	\item Isola fino a 1500 volt -Alimentazione da+/-4,5V a +/-18V.
	\item Possiede 2 alimentazioni separate e quindi anche due riferimenti di massa separati.
	\item Non richiede componenti esterni (tranne le capacità di disaccoppiamento)
\end{itemize}


\begin{figure}[tpc]
\centering
\includegraphics[width=.7\textwidth]{../E07/latex/circuito.pdf}
\caption{Schema del circuito da noi utilizzato per acquisire l'elettrocardiogramma.}
\label{cir8:compensation}
\end{figure}

%\subsection{SECONDO PARAGRAFO MEGLIO DEL PRIMO E DELL'INTRODUZIONE DEL GRUPPO AVANZATO}
\subsection*{Conclusioni}
In questa esperienza siamo riusciti ad ottenere un elettrocardiogramma affetto da poco rumore, grazie agli stadi di abbattimento inseriti, e amplificato in maniera tale da poter essere letto facilmente dalla strumentazione.


\clearpage
\input{relazione8.tex}
\clearpage
\section{11.11.2014 - Porte logiche}

In questa esperienza verificheremo il funzionamento di alune porte logiche TTL e semplici circuiti costruiti con esse.

\subsection*{Strumenti e materiali}

%\begin{figure}[htc]
\begin{itemize} [noitemsep]
	%\item Generatore di forme d'onta Agilent 33120A con range di frequenza da \SI{100}{\micro\hertz} a \SI{15}{\mega\hertz};
	\item Oscilloscopio Agilent DSO-X 2002A (bandwidth \SI{70}{\mega\hertz}, sample rate \num{2} GSa/s);%\newline
%	\begin{minipage}{0.65\textwidth}
%		\vspace{0.4mm}
%%		\begin{itemize} [noitemsep]
%%		\item Oscilloscopio Agilent DSO-X 2002A (bandwidth \SI{70}{\mega\hertz}, sample rate \num{2} GSa/s);
		\item Generatore di tensione continua Agilent E3631A (max $\pm \, \SI{25}{\volt}$ o $\pm \, \SI{6}{\volt}$);
%%		\item Generatore di forme d'onta Agilent 33120A con range di frequenza da \SI{100}{\micro\hertz} a \SI{15}{\mega\hertz};
		\item Multimetro Agilent 34410A a sei cifre e mezza;
		\item Amplificatori operazionali OP07;
		\item Un integrato AD622 e un ISO124;
		\item Un integrato 7400;
		\item Basetta a LED;
		%\item Transistor di potenza NPN 2N2222;
		\item Porte logiche TTL;		
		\item Resistenze e capacità di vari valori;
		%\item Trimmer multigiro da \SI{10}{\kilo\ohm} e \SI{1}{\kilo\ohm};
		\item Breadboard e cablaggi vari.
%%		\end{itemize}
%	\end{minipage}
%	\begin{minipage}{0.3\textwidth}
%			\centering
%			\includegraphics[width=3.5cm]{../E06/latex/REF02.pdf}
%			\caption{Piedinatura dell'integrato\newline REF02.}
%			\label{cir6:REF02}
%	\end{minipage}
\end{itemize}
%\end{figure}
%\vspace{-.8cm}


\subsection{Porte TTL e scheda di visualizzazione}
Le porte logiche a tecnologia TTL funzionano utilizzando come 1 logico il valore di tensione +\SI{5}{\volt} mentre come 0 logico \SI{0}{\volt}. Per controllare il funzionamento dei nostri circuiti utilizzeremo una schedina di visualizzazione a led. Tale circuito permette di visualizzare lo stato "alto" o "basso" delle linee da monitorare tramite 8 led (4 rossi e 4 verdi). Tali schede contengono un integrato buffer 74LS244 nella logica TTL che fornisce in uscita una corrente di \SI{-15}{\milli\ampere} per il livello alto e di \SI{24}{\milli\ampere} per il livello basso.
Ricordiamo che l'alimentazione deve essere stabilizzata e disaccoppiata con i condensatori come nel caso analogico in quanto nei momenti di commutazione delle porte abbiamo grandi variazioni di tensione in piccoli tempi (idealmente la variazione da stato 1 a 0 è istantanea). È dunque essenziale avere una riserva di cariche utilizzabili dal circuito nel momento in cui si hanno commutazioni.

\subsubsection{Porta NAND}


È stato montato il circuito come 


\begin{tabular}{|l|l|l|}
\hline
A & B & $\overline{AB}$ \\
\hline
0 & 0 & 1\\
\hline
0 & 1 & 1\\
\hline
1 & 0 & 1\\
\hline
1 & 1 & 0\\
\hline
\end{tabular}

\subsubsection{Porta NOT}


\begin{tabular}{|l|l|||l|l|}
\hline
$V_{in}$ [\si{\volt}] & $V_{out}$ [\si{\volt}] & $V_{in}$ [\si{\volt}]& $V_{out}$ [\si{\volt}]\\
\hline
0.5 &  & 3 &  \\
\hline
1 &  & 3.5 &  \\
\hline
1.5 &  & 4 &  \\
\hline
2 &  & 4.5 &  \\
\hline
2.5 &  & 5 &  \\
\hline
\end{tabular}


\subsubsection{Circuito di GATE}

\subsubsection{Porta XOR}

\subsubsection{Votazione con 3 Giurati e 1 Presidente}

\subsubsection{Allarme MINI-APPARTAMENTO}






%\subsection{Porte TTL e scheda di visualizzazione}
Le porte logiche a tecnologia TTL funzionano utilizzando come 1 logico il valore di tensione +\SI{5}{\volt} mentre come 0 logico \SI{0}{\volt}. Per controllare il funzionamento dei nostri circuiti utilizzeremo una schedina di visualizzazione a led. Tale circuito permette di visualizzare lo stato "alto" o "basso" delle linee da monitorare tramite 8 led (4 rossi e 4 verdi). Tali schede contengono un integrato buffer 74LS244 nella logica TTL che fornisce in uscita una corrente di \SI{-15}{\milli\ampere} per il livello alto e di \SI{24}{\milli\ampere} per il livello basso.
Ricordiamo che l'alimentazione deve essere stabilizzata e disaccoppiata con i condensatori come nel caso analogico in quanto nei momenti di commutazione delle porte abbiamo grandi variazioni di tensione in piccoli tempi (idealmente la variazione da stato 1 a 0 è istantanea). È dunque essenziale avere una riserva di cariche utilizzabili dal circuito nel momento in cui si hanno commutazioni.

\subsubsection{Porta NAND}

\begin{wrapfigure}[5]{r}{0.45\textwidth}
\centering
\includegraphics[width=.15\textwidth]{../E09/latex/NAND.pdf}
\caption{Simbolo circuitale convenzionalmente utilizzato per porte NAND.}
\label{cir9:nand}
\end{wrapfigure}

Abbiamo verificato il funzionamento di una porta NAND utilizzando, come già detto, una scheda di visualizzazione a LED. È stata verificata la seguente tavola di verità. 



\begin{tabular}{|l|l|l|}
\hline
A & B & $\overline{AB}$ \\
\hline
0 & 0 & 1\\
\hline
0 & 1 & 1\\
\hline
1 & 0 & 1\\
\hline
1 & 1 & 0\\
\hline
\end{tabular}

\subsubsection{Porta NOT}


\begin{wrapfigure}[10]{r}{0.35\textwidth}
\centering
\includegraphics[width=.15\textwidth]{../E09/latex/NOT.pdf}
\caption{Implementazione e simbolo circuitale per porte NOT.}
\label{cir9:not}
\end{wrapfigure}


Non avendo a disposizione una porta NOT, abbiamo utilizzato una porta NAND cortocircuitando i due ingressi tra loro. Così facendo otteniamo un'implementazione di una porta NOT. Ne abbiamo in primis verificato il funzionamento collengandola alla scheda di visualizzazione led e fornendo un'onda quadra in ingresso ($V_{pp}=\SI{5}{\volt}$, $V_{off}=+\SI{2.5}{\volt}$, $\nu=\SI{2}{\hertz}$). Una volta verificato funzionasse, collegato ingresso ed uscita all'oscilloscopio e ne abbiamo analizzato la risposta per diverse tensioni costanti in ingresso. I dati ottenuti sono riportati in Tabella (\ref{tab9:risposta}). Come vediamo, tra \SI{1}{\volt} e \SI{1.5}{\volt} la tensione varia in modo molto rapido.

Siamo dunque interessati a capire meglio la risposta della porta in funzione della tensione in ingresso. Per fare ciò ci aiuta l'oscilloscopio con la modalità XY.  




\begin{tabular}{|l|l|||l|l|||l|l|||l|l|}
\hline
$V_{in}$ [\si{\volt}] & $V_{out}$ [\si{\volt}] & $V_{in}$ [\si{\volt}]& $V_{out}$ [\si{\volt}] & $V_{in}$ [\si{\volt}]& $V_{out}$ [\si{\volt}] & $V_{in}$ [\si{\volt}]& $V_{out}$ [\si{\volt}]\\
\hline
0 & 4.30 & 2.5 & 0.10 & 5  & 0.10 & 2.5 & 0.10\\
\hline
0.5 & 4.05 & 3 & 0.10 & 4.5 & 0.10 & 2 &0.10\\
\hline
1 & 3.01 & 3.5 & 0.10 & 4 & 0.10 & 1.5 & 0.10\\
\hline
1.5 & 0.10 & 4 & 0.10 & 3.5 & 0.10 & 1 & 3.00\\
\hline
2 & 0.10 & 4.5 & 0.10 & 3.0 &0.10 & 0.5 & 4.05\\
\hline
\label{tab9:risposta}
\end{tabular}


Abbiamo dunque fornito un'onda sinusoidale all'ingresso di frequenza \SI{100}{\kilo\hertz}, $V_{pp}=\SI{5}{\volt}$ e $V_{off}=+\SI{2.5}{\volt}$. Come vediamo dal grafico sotto riportato, sebbene stiamo lavorando con segnali digitali (1 o 0) le tensioni in gioco sono manifestamente analogiche. Avremo dunque tutta la serie di problemi già affrontati nella parte analogica del corso. Vediamo infatti che la tensione in uscita non è una funzione a gradino. Inoltre, per evitare che rumori molesti disturbino il segnale in uscita, avremo bisogno di una sorta di isteresi. Non esisterà dunque un valore ben definito di V alta e V bassa, ma delle bande di lavoro. Solitamente nel data-shit sono riportati tali parametri.

$$GRAFICO$$

\subsubsection{Circuito di GATE con porta AND}
Per realizzare il circuito GATE è sufficiente utilizzare una porta AND con un segnale di controllo ad un ingresso. Per realizzare una porta AND possiamo utilizzare 2 NAND come proposto nella seguente figura. Infatti, utilizzando la notazione convenzionale, $AB=\overline {\overline {AB}}$ che trivialmente è una NAND e una NOT in serie.
 
\begin{wrapfigure}[10]{r}{0.35\textwidth}
\centering
\includegraphics[width=.15\textwidth]{../E09/latex/AND.pdf}
\caption{Implementazione e simbolo circuitale per porte AND.}
\label{cir9:and}
\end{wrapfigure}

La tavola di verità è la seguente (C=controllo, S=segnale):

\begin{tabular}{|l|l|l|}
\hline
C & S & CS \\
\hline
0 & 0 & 0\\
\hline
0 & 1 & 0\\
\hline
1 & 0 & 0\\
\hline
1 & 1 & 1\\
\hline
\end{tabular}

Come vediamo dalla tabella, se il segnale di controllo è a 0 logico l'uscita sarà sempre 0 logico. Se invece il controllo è a 1 logico, allora l'uscita sarà una copia del segnale in ingresso. Per verificare ciò, abbiamo collegato l'oscilloscopio all'uscita del circuito e fornito un'onda quadra ($V_{pp}=\SI{5}{\volt}$, $V_{off}=+\SI{2.5}{\volt}$, $\nu=\SI{2}{\hertz}$). Abbiamo visto che quando il controllolol era collegato a comune, l'uscita era praticamente nulla. Nell'altro caso invece avevamo la stessa onda quadra dell'ingresso anche all'uscita.

\subsubsection{Porta XOR}

Una porta XOR implementa la funzione matematica della somma diretta. Avendo solo a disposizione porte NAND, dobbiamo cercare di implementare un circuito che svolga tale funzione. Riportiamo in tabella la tavola di verità. In questo semplice caso non è necessario scomodare dalla tomba il povero Karnaugh.

\begin{tabular}{|l|l|l|}
\hline
A & B & $A \oplus B$ \\
\hline
0 & 0 & 0\\
\hline
0 & 1 & 1\\
\hline
1 & 0 & 1\\
\hline
1 & 1 & 0\\
\hline
\end{tabular}

\begin{wrapfigure}[15]{r}{0.35\textwidth}
\centering
\includegraphics[width=.30\textwidth]{../E09/latex/XOR.pdf}
\caption{Implementazione e simbolo circuitale per porte XOR.}
\label{cir9:xor}
\end{wrapfigure}

Dunque $A \oplus B=A\overline B + \overline A B$, dove con A,B indichiamo l'1 logico, con $\overline A,\overline B$ lo 0 logico.

Applichiamo le regole di somme e prodotti viste a lezione:

$$A \oplus B=A\overline B + \overline A B=\overline{\overline{A\overline B + \overline A B}}=\overline{\overline{A\overline B} \cdot \overline{\overline A B}}$$

Dunque, il circuito per realizzare uno XOR usando porte NAND e NOT è il seguente.



Come vediamo, la realizzazione da noi progettata prevede l'utilizza do 2 porte NOT e 3 NAND. Ne abbiamo verificato il funzionamento con la solita scheda a LED.





\subsubsection{Votazione con 3 Giurati e 1 Presidente}
Il circuito di votazione di 3 Giurati con 1 Presidente prevede di sommare i voti dei vari membri tenendo però in considerazione che il voto del presidente vale doppio. Partiamo dunque dalla mappa di Karnaugh.


\begin{tabular}{|l|l|l|l|l|}
\hline
\diaghead{\theadfont lololololo a} {CP}{AB}& 00& 01 & 11&10\\
\hline
00&0 & 0 & 0 &0 \\
\hline
01&0 & 1 & 1 &1 \\
\hline
11&1 &1  &1  &1 \\
\hline
10&0 & 0 & 1 & 0\\

\hline

\end{tabular}




Ora, minimizzando tale mappa otteniamo $Y=ABC+P(A+B+C)$. Come fatto prima, avendo a disposizione solo porte NAND, trasformiamo tutto in prodotti.

$$Y=ABC+P(A+B+C)=\overline{\overline{ABC+P(A+B+C)}}=\overline{\overline{ABC} \cdot \overline {P(A+B+C)} }= \overline{\overline{(AB)C} \cdot \overline {P(\overline{\overline{{A+B+C}} })}}=\overline{\overline{(AB)C} \cdot \overline {P(\overline{{(\overline A \cdot \overline B) \cdot \overline C} })}}$$

Ora è immediato costruire il circuito utilizzando porte NAND e NOT. Lo schema è riportato in figura. Come per gli altri circuiti, è stato verificato utilizzando la schedina a LED. 



\begin{figure}[htpc]
\centering
\includegraphics[width=.75\textwidth]{../E09/latex/giudici.pdf}
\caption{Buzz non ha capito che sono 3 giurati e 1 giudice, non 4 giudici.}
\label{cir9:giudici}
\end{figure}


\subsubsection{Allarme MINI-APPARTAMENTO}

In quest'ultima parte dell'esperienza cercheremo di progettare e costruire un circuito di allarme. Immaginiamo di avere una villa piena di zoccole e di aver installato dei sensori. Vogliamo analizzarne i segnali per decidere se l'allarme deve suonare o no. Chiamiamo P il sensore sulla porta (0=chiusa), F il sensore sulla finestra (0=chiusa), I il sensore ad infrarossi (0=no movimento) e infine C la chiave che permette di attivare l'allarme escludendo il sensore ad infrarossi (1=escludo infrarossi). 

Costruiamoci la mappa di Karnaugh.


\begin{tabular}{|l|l|l|l|l|}
\hline
\diaghead{\theadfont lololololo a} {IC}{PF}& 00& 01 & 11&10\\
\hline
00&0 & 1 & 1 &1 \\
\hline
01&0 & 1 & 1 &1 \\
\hline
11&0 &1  &1  &1 \\
\hline
10&1 & 1 & 1 & 1\\

\hline

\end{tabular}

Trivialmente, una dei possibili raggruppamenti è il seguente $Y=P+F+\overline C I$. Con la ormai nota tecnica trasformiamo tutto in NAND e NOT.

$$Y=P+F+\overline C I=\overline{\overline{P+F+\overline C I}}=\overline{(\overline P \cdot \overline F) \cdot \overline{\overline C I}}$$ 

Il circuito risulta quello nella seguente figura. È stato verificato come gli altri utilizzando la scheda a LED.



\begin{figure}[htpc]
\centering
\includegraphics[width=.65\textwidth]{../E09/latex/allarme.pdf}
\caption{Bisogna negare il C, non I.}
\label{cir9:allarme}
\end{figure}








%\subsection{SECONDO PARAGRAFO MEGLIO DEL PRIMO E DELL'INTRODUZIONE DEL GRUPPO AVANZATO}
%\subsection{Verifica della banda passante}

In questa parte dell'esperienza vogliamo valutare la banda passante di un circuito amplificatore (non invertente)
%\subsection{Guadagno Open Loop}
\label{par3:open}

\begin{wrapfigure}[17]{r}{0.45\textwidth}
  \begin{center}
    \includegraphics[width=0.280\textwidth]{../E03/latex/LF_ol.pdf}
  \end{center}
  \caption{Circuito utilizzato per misurare il guadagno $A_{ol}$. I valori di $R_4$ utilizzati sono $(5.2\pm 0.1)$ \si{\ohm},  $(9.6\pm0.1)$ \si{\ohm}, $(101.3\pm 0.1)$ \si{\ohm} e $(998.2\pm 0.1)$ \si{\ohm} mentre $R_1=\SI{98\pm1}{\kohm}$, $R_2=\SI{97\pm1}{\kohm}$ e $R_3=\SI{98\pm1}{\kohm}$.}
  \label{cir3:low_frequency}
\end{wrapfigure}

In questa ultima parte dell'esperienza abbiamo cercato di misurare il guadagno $A_{ol}$ del nostro amplificatore reale.
Come sappiamo, tale guadagno è una funzione della frequenza e dobbiamo dunque trovare un modo per misurarlo.

Fino alla frequenza di circa \SI{8}{\hertz} il guadagno di un opamp $\mu$A741 è costante a circa \num{2E5}.
Oltre tale frequenza abbiamo una caduta di circa \SI{20}{\decibel}/decade.
Non risulta dunque possibile effettuare misure a loop aperto in quanto piccolissime variazioni di tensione ai due ingressi causerebbero grandi effetti in uscita.

Abbiamo dunque progettato un circuito per controllare la tensione in uscita così da non mandare in saturazione il nostro opamp.
In Figura \ref{cir3:low_frequency} è riportato lo schema circuitale da noi utilizzato.
Utilizzando l'oscilloscopio abbiamo misurato i valori di tensione $V_A$ e $V_{out}$. Cerchiamo dunque di legare il guadagno a tali quantità.

Nel caso dell'amplificatore operazionale sappiamo che vale l'eq. (\ref{eq3:regola_opamp}), dove $(V_+-V_-)$ è semplicemente la differenza di potenziale presente tra i due ingressi, che sarà ovviamente data da $\Delta V = I_3R_4$, dove $I_3$ è la corrente che scorre attraverso $R_3$ ed $R_4$\footnote{Assumiamo che la corrente assorbita dagli ingressi sia trascurabile.}.
È quindi facile stimare $I_3=\frac{V_A}{R_3+R_4}$ conoscendo $V_A$.

Si ottiene dunque:
\begin{equation}
A_{ol}=\frac{{V_{out}}^{pp}}{{V_A}^{pp}} \frac{R_4+R_3}{R_4}
\label{eq3:lfgain}
\end{equation}

Osserviamo come in equazione (\ref{eq3:lfgain}) non compaia il termine $V_{in}$.
Dovremo dunque scegliere per le varie frequenze una tensione picco-picco in ingresso adeguata in modo che il nostro opamp non saturi e allo stesso tempo $V_A$ sia sufficientemente grande da essere poco influenzata dal rumore di fondo e sia stabile il più possibile.

Durante le misurazioni abbiamo deciso di variare la resistenza $R_4$ suddividendo il procedimento di misura in diversi step; questa strategia ci ha permesso di mantenere i valori di $V_A$ e $V_{out}$ sempre nell'ordine del volt. 

\subsubsection*{Misura per alte frequenze}

\begin{wrapfigure}[13]{r}{0.45\textwidth}
  \begin{center}
    \includegraphics[width=0.28\textwidth]{../E03/latex/HF_ol.pdf}
  \end{center}
  \caption{Circuito utilizzato per stimare il guadagno $A_{ol}$ a basse frequenze. I valori delle resistenze utilizzate sono le stesse del circuito in Figura \ref{cir3:low_frequency}.}
  \label{cir3:high_frequency}
\end{wrapfigure}

Per non dover utilizzare resistenze $R_4$ troppo grandi ad alte frequenze ($>$ \SI{70}{\kilo\hertz}), abbiamo provato ad effettuare le misure direttamente ad open-loop, dato che il guadagno è contenuto (\num{<100}).
Tuttavia il generatore di forme d'onda non emette un segnale pulito alla sola frequenza voluta, in quanto anch'esso è uno strumento reale e dunque affetto da errori.
Esso potrebbe dunque emettere segnali a frequenze basse o addirittura un piccolo offset.
Poichè per tali frequenze il guadagno dell'amplificatore operazionale è enorme, il segnale in uscita viene immediatamente distorto da tali rumori indesiderati.
Non è da sottovalutare poi la presenza dei residui dell'offset dell'opamp.
Esso infatti dipende anche dalla temperatura e dunque la sua compensazione non potrà mai essere perfetta ad ogni temperatura.
Misurare il guadagno $A_{ol}$ effettivamente a open-loop è particolarmente problematico.

Per aggirare tali ostacoli possiamo utilizzare il circuito riportato in Figura \ref{cir3:low_frequency}, rimuovendo la resistenza $R_4$ e sostituendo $R_3$ con un filo partendo dalla configurazione di Figura \ref{cir3:high_frequency}.
Cosi facendo otteniamo la semplice relazione $A_{ol}=\frac{{V_{out}}^{pp}}{{V_A}^{pp}}$, che possiamo anche vedere come limite di eq. (\ref{eq3:lfgain}) per $R_4 \rightarrow \infty$.
I dati ottenuti sono riportati nel seguente grafico. 

\begin{figure}[h]
	\centering
	\includegraphics[width=0.9\textwidth]{../E03/latex/gol.pdf}
	\caption{Grafico semilogaritmico del guadagno (in $dB$) in funzione della frequenza. Ad ogni forma dei punti (e relativo colore) corrisponde una diversa resistenza $R_4$ utilizzata. Per quanto riguarda $R_4=\infty$, intendiamo la mancanza del collegamento fra ingresso invertente e non invertente (circuito in Figura \ref{cir3:high_frequency}). La legge plottata è la (\ref{eq3:fit_gain_ol}).}
  \label{cir3:gain_open_loop}
\end{figure}

Anche in questo caso, considerando l'operazionale come un semplice passa basso, possiamo utilizzare l'equazione (\ref{eq3:passa_basso}) e trovare la relazione fra il guadagno e la frequenza. Otteniamo
\begin{equation}
\mathrm{Gain}=\frac{|V_{out}|}{|V_{in}|}=\sqrt{\frac{A_{ol}^2 \omega_0^2}{\omega^2 + \omega_0^2}}
\label{eq3:fit_gain_ol}
\end{equation}

Per un fit e i parametri si veda il paragrafo successivo.

\subsection{Calcolo di $A_{ol}$}
\label{par3:A_ol}

\subsubsection{$A_{ol}$ da circuiti retroazionati}

Con un fit sulle leggi (\ref{eq3:fit_gain}) e (\ref{eq3:fit_fase}) abbiamo ottenuto una stima dei parametri di $A_{ol}$ e $\omega_0$. Abbiamo deciso, dato che l'opamp utilizzato era lo stesso per tutte le misurazioni, di procedere con una media pesata sui valori ottenuti con i 4 fit (due con (\ref{eq3:fit_gain}) e due con (\ref{eq3:fit_fase})), ottenendo i seguenti valori sperimentali
$$A_{ol}=(1.8\pm0.1)\times 10^5 \qquad f_0=(5.9\pm0.2)\si{\hertz}$$

Ovviamente possiamo ottenere un valore di $A_{ol}$ anche utilizzando la (\ref{eq3:freq_taglio})
$$A_{ol} = \frac{G}{\omega_0} (\omega_t - \omega_0)$$
e considerando la frequenza $f_0$ ricavata dal fit. Si ottiene un valore teorico (media dei valori con i due guadagni del circuito) di $A_{ol} = (1.71\pm 0.06)\times 10^5$, compatibile con il valore sperimentale.

\subsubsection{$A_{ol}$ da open loop}

Con il fit dalla (\ref{eq3:fit_gain_ol}) otteniamo i seguenti valori
$$A_{ol}=(2.0\pm0.1)\times 10^5 \qquad f_0=(5.8\pm0.2)\si{\hertz}$$

I dati ottenuti non risultano compatibili ma riteniamo tale fatto sia dovuto ad errori parassiti non ben identificati.
\subsection*{Conclusioni}
In questa esperienza abbiamo verificato il funzionamento di alcune porte, evidenziando il fatto che sebbene i segnali trattati siano digitali (0 e 1 logico, 0 e $+$\SI{5}{\V}), non possiamo trascurare la tecnologia analogica che sta alla base del funzionamento.
Inoltre, abbiamo visto come con l'utilizzo delle più comuni porte NAND, possiamo implementare le altre porte.
Infine, utilizzando le proprierà algebriche della logica booleana, abbiamo semplificato circuiti complessi e li abbiamo implementati con il solo utilizzo di porte NAND.




\end{document}
